\documentclass[12pt,oneside,a4paper]{book}
\usepackage[latin1]{inputenc}
\usepackage[english, spanish]{babel}	%For languages characters and hyphenation
\usepackage{amsmath}
\usepackage{amsfonts}
\usepackage{titlesec}
\usepackage{amssymb}
\usepackage{graphicx}
\usepackage{array}
\usepackage{cite}
\usepackage[all]{xy}
\usepackage{fancyhdr}
\usepackage{lastpage}
\usepackage{caption}
\usepackage{subcaption}
\usepackage{multicol}
\usepackage{tabularx}
\usepackage{hhline}
\usepackage{multirow}
\usepackage{microtype}		%Makes pdf look better.
\usepackage[left=3.00cm, right=3.00cm, top=2.00cm, bottom=2.00cm]{geometry}
\author{Julian David Mora Ramos}
\title{Proyecto, implementaci�n de SIGEPI}
\usepackage[pdftex,
	pdfauthor={JULIAN DAVID MR},
	pdftitle={IMPLEMENTACI�N DEL SISTEMA SIGEPI},
	pdfsubject={DESARROLLO DE SOFTWARE},
	pdfkeywords={SOFTWARE, MDD, MDE, UDLA, DOCUMENTACI�N, DESARROLLO},
	pdfproducer={Universidad de la Amazonia},
	pdfcreator={@anlijudavid}]{hyperref}

% -------------------------
\setcounter{tocdepth}{4}
\setcounter{secnumdepth}{4}
\newcommand{\myparagraph}[1]{\paragraph{#1}\mbox{}\\}	
% -------------------------

% ----------------Tama�o de texto Chapters-----------------
\newcommand{\chapfnt}{\fontsize{16}{19}}
\newcommand{\secfnt}{\fontsize{14}{17}}
\newcommand{\ssecfnt}{\fontsize{12}{14}}

\titleformat{\chapter}[display]
{\normalfont\chapfnt\bfseries}{\chaptertitlename\ \thechapter}{20pt}{\chapfnt}

\titleformat{\section}
{\normalfont\secfnt\bfseries}{\thesection}{1em}{}

\titleformat{\subsection}
{\normalfont\ssecfnt\bfseries}{\thesubsection}{1em}{}

\titlespacing*{\chapter} {0pt}{50pt}{40pt}
\titlespacing*{\section} {0pt}{3.5ex plus 1ex minus .2ex}{2.3ex plus .2ex}
\titlespacing*{\subsection} {0pt}{3.25ex plus 1ex minus .2ex}{1.5ex plus .2ex}
% --------------------------------------------------------

% -----------------Config de grafica MVC------------------
\usepackage{tikz}
\usetikzlibrary{arrows,positioning} 
\tikzset{
	%Define standard arrow tip
	>=stealth',
	%Define style for boxes
	punkt/.style={
		rectangle,
		rounded corners,
		draw=black, very thick,
		text width=6.5em,
		minimum height=2em,
		text centered},
	% Define arrow style
	pil/.style={
		->,
		thick,
		shorten <=2pt,
		shorten >=2pt,}
}
% --------------------------------------------------------

% ------------Permite renderizar muchas tablas------------
\maxdeadcycles=1000
% --------------------------------------------------------

\usepackage{float} % para usar [H]

% --------------------------------------------------------
% --------------------------------------------------------


\begin{document}
	\markboth{"SIGEPI"}{''Universidad de la Amazonia''}
	
	% Paginas iniciales
	\title{%
	\begin{figure}[htbp]
		\centering
		\includegraphics[width=50px]{resources/logoudla.jpg}
	\end{figure} \\ %
	Implementaci�n del sistema de informaci�n SIGEPI basado en herramientas de desarrollo dirigido por modelos \\
	\large Implementaci�n de la plataforma web para la gesti�n de los proyectos presentados por los grupos y semilleros de investigaci�n en la Universidad de la Amazonia \\
}
\author{
	Julian David Mora Ramos $^{1}$
	\\
	\small{$^{1}$Universidad de la Amazonia, $^{1}$Florencia,Caquet�}\\
}
\maketitle
	\thispagestyle{empty}

\begin{center}
	\includegraphics{./resources/logoudla.jpg}
\end{center}
\vspace{2cm}

\begin{center}
	{\Large {\bf UNIVERSIDAD DE LA AMAZONIA}}
	\vspace{5mm}

 	{\Large {Ingenier�a de Sistemas}}
    \\
  	\vspace{34mm}
	{\large {\bf PROYECTO FIN DE PREGRADO}}
  	\vspace{10mm}
    \\
  	{\Large {{\Huge {				
  		IMPLEMENTACI�N DEL SISTEMA INFORMACI�N PARA GESTI�N DOCUMENTAL SIGEPI
	}} \\[1cm] }}
  	\vspace{2cm}
	{\large {
  		{\bf Autor}:\\
  		Julian David Mora Ramos\\
  	}}
	\vspace{10mm}
  	{\large {Curso acad�mico 2011/2017}}
  	\vspace{1cm}
\end{center}
	%------------------------------------------------
%   DEDICATORIA
%------------------------------------------------

\pagestyle{empty}
\frontmatter

\chapter*{}
\begin{flushright}
	\textit{DEDICATORIA}
		
	A la Facultad de Ingenier�a y a la  Universidad, por la formaci�n que me han dado.\\
	Es gracias a ustedes que es posible el presente trabajo.\\
	En verdad, gracias.\\
	Yo.
\end{flushright}
	%------------------------------------------------------
%   AGRADECIMIENTOS
%------------------------------------------------------

\chapter*{Agradecimientos}
%\markboth{AGRADECIMIENTOS23}{AGRADECIMIENTOS} % encabezado 

�Muchas gracias a todos!

	\section*{Nota de aceptaci�n}
	
Tambi�n quisiera reconocer a ... por ...CONACYT,  PAPIIT / etc.



	%--------------------------------------------------
%   DECLARACI�N
%--------------------------------------------------

\thispagestyle{empty}
\vspace*{\fill}
\begingroup
``Con fundamento en los art�culos 21 y 27 de la Ley Federal del Derecho de Autor y como titular de los derechos moral y patrimonial de la obra titulada ``\textbf{T�TULO DE LA TESIS}'', otorgo de manera gratuita y permanente al Instituto Tecnol�gico Aut�nomo de M�xico y a la Biblioteca Ra�l Baill�res Jr., la autorizaci�n para que fijen la obra en cualquier medio, incluido el electr�nico, y la divulguen entre sus usuarios, profesores, estudiantes o terceras personas, sin que pueda percibir por tal divulgaci�n una contraprestaci�n''.

\centering

\hspace{3em}

\textsc{AUTOR}

\vspace{5em}

\rule[1em]{20em}{0.5pt} % L�nea para la fecha

\textsc{Fecha}
 
\vspace{8em}

\rule[1em]{20em}{0.5pt} % L�nea para la firma

\textsc{Firma}

\endgroup
\vspace*{\fill}

	\begin{center}	
	\section*{Resumen}
	
	
	\section*{Abstract}
\end{center}
	
	\tableofcontents
	
	\cleardoublepage
	\addcontentsline{toc}{chapter}{Lista de figuras} % para que aparezca en el indice de contenidos
	\listoffigures % indice de figuras
	
	\cleardoublepage
	\addcontentsline{toc}{chapter}{Lista de tablas} % para que aparezca en el indice de contenidos
	\listoftables % indice de tablas
	
	% \pagenumbering{arabic} % para empezar la numeraci�n con n�meros
	% Empieza la numeraci�n de las p�ginas	
	\mainmatter
		
	% Contenido
	\section[INTRODUCCI�N]{INTRODUCCI�N}

Actualmente, adaptarse a las necesidades del cliente es un
problema que aumenta con el transcurso del tiempo, las condiciones
de negocio cambian y consigo tambi�n la construcci�n de sistemas.
Consecutivamente, el mundo moderno exige gran demanda tecnol�gica
y adaptarse a esta linea de constantes transformaciones resulta
en una tarea compleja. Por esto, gracias a las investigaciones
realizadas en este campo, se han logrado avances en ciertas �reas
fundamentales que influyen en la puesta en marcha de proyectos.
\\

Muchos proyectos de desarrollo de software poseen gran demanda, y
m�s cuando est�n soportados por herramientas de automatizaci�n
de tareas. Los tiempos de ejecuci�n y elaboraci�n de proyectos se
reducen cuando herramientas DSL (Lenguaje de dominio especifico)
son implementadas, ofreciendo caracter�sticas de generaci�n de
c�digo reutilizable y componentes.
\\

El lenguaje de dominio especifico no es una tecnolog�a emergente,
las primeras apariciones fueron a mediados de los 80s. 
Actualmente se pueden encontrar ramificaciones DSL 
en diferentes aplicaciones, cada una con
funcionalidades especificas, tales como MDD (Desarrollo Dirigido
por Modelos), MDA (Arquitectura Dirigida por Modelos) y
MDE (Ingenieria Orientada a Modelos), siendo estos un conjunto de
recursos interoperables que permiten ser usados en el
an�lisis, modelado y construcci�n de proyectos
software. Por otro lado, han aumentado la cantidad de mecanismos
para agilizar la planeaci�n y ejecuci�n de proyectos, generalmente
todos estos siguen la misma filosof�a de las DSL, por ejemplo, se
pueden encontrar las ORM (Mapeo de Objetos Relacionales) para el
control de m�ltiples gestores
de bases de datos a partir de un solo lenguaje base. Del mismo 
modo, los ingenieros, programadores de computadoras, matem�ticos,
estad�sticos y dem�s, hacen uso de herramientas de lenguaje de
dominio 
especifico a partir de programas de computador como R (Lenguaje
estad�stico), interpretes 
de expresiones regulares, SQL (lenguaje de consulta estructurada), 
LINQ (Lenguaje de consulta integrada), entre otros. Como se puede
observar, un dominio espec�fico tiene un vocabulario especializado
para describir las cosas que son particulares a ese dominio.
\\

El sistema de gesti�n documental, SIGEPI, es una plataforma web 
robusta, creada gracias a la implementaci�n de herramientas 
de modelado y generaci�n de c�digo, DSL. En la primera secci�n de 
este documento se nombran todos aquellos procesos que se siguieron 
para la obtenci�n del c�digo fuente a partir de los esquemas 
dados al aplicar MDD (Desarrollo Dirigido por Modelos).
Posteriormente, se explica en detalle la arquitectura de desarrollo que fue
aplicada, incluyendo cada una de las librer�as y marcos de trabajo que 
fueron usados.
	\section[DESCRIPCI�N GENERAL DEL PROYECTO]{DESCRIPCI�N GENERAL DEL PROYECTO}

\subsection[Planteamiento del problema]{Planteamiento del problema}

\subsubsection[Contexto]{Contexto}
\subsubsection[Formulaci�n del problema]{Formulaci�n del problema}

\subsection[Justificaci�n]{Justificaci�n}

\subsection[Objetivos]{Objetivos}
\subsubsection[Objetivo general]{Objetivo general}
\subsubsection[Objetivo especifico]{Objetivo especifico}

	\section[MARCO REFERENCIAL]{MARCO REFERENCIAL}

\subsection[Marco te�rico]{Marco te�rico}

\subsubsection[Desarrollo Dirigido por Modelos. Conceptos]{Desarrollo Dirigido por Modelos. Conceptos}
El desarrollo dirigido por modelos (MDD) es un paradigma que resuelve
muchos inconvenientes en el desarrollo de software, la causa 
surgi� desde los inicios de la d�cada de los 60s, cuando se
introdujo el concepto de ''crisis del software'', originado por la complejidad y el costo requerido por las necesidades del cliente \cite{pons2010desarrollo}.
MDD enriquece el proceso de desarrollo a partir de herramientas 
especializadas, destaca una relaci�n entre lo abstracto y el c�digo
fuente. La abstracci�n est� representada por modelos, por tanto 
estos requieren de una o m�s fases de transformaci�n para 
finalmente obtener el c�digo [ver figura \ref{fig_transcode}].

\begin{displaymath}
	\xymatrix{
		*+<1cm>[F-,]{Modelo}\ar[r] 
		& *+<1cm>[o][F]{M2T/M2M}\ar[r]
		& ... \ar[r]
		& *+<1cm>[F-,]{Modelo}\ar[r]
		& *+<.5cm>[F]{Codigo}
	}
\end{displaymath}
\captionof{figure}{Transformaci�n de c�digo}
\label{fig_transcode}


\subsubsection[DSL, Lenguaje de dominio especifico]{Lenguaje de dominio especifico}
El lenguaje de dominio especifico eleva el nivel de abstracci�n m�s all�
que los lenguajes de programaci�n para especificar una soluci�n al
problema usando conceptos de dominio \cite{kelly2008domain}. 
Generalmente se le encuentra como una notaci�n gr�fica, cuyos
modelos resultan en un conjunto de elementos y relaciones entre si.

\begin{displaymath}
	\xymatrix{
		\underline{Requerimientos}\ar@/_1mm/[r] 
		& *+<1cm>[o][F]{\texttt{Abstracci�n}}\ar[r]
		& *+<1cm>[F-,]{\textbf{DSL}}
	}
\end{displaymath}
\captionof{figure}{Abstracci�n de DSL}
\label{fig_abs_dsl}

\subsubsection[Meta modelos]{Meta modelos}


\subsubsection[Arquitectura MVC]{Arquitectura MVC}

\subsubsection[Framework de desarrollo Ext.NET]{Framework de desarrollo Ext.NET}
	\chapter{Procesos para presentaci\'on de proyectos}

Vicerector\'ia de investigaciones de la Universidad de la
Amazonia cuenta 

	\chapter{Base de datos}

Diccionario de datos 

\begin{table}[ht]
	\caption{Tabla Usuario}
	\label{labelTableUsuario}
	\begin{tabular}{ |l|l|l|l|l|l| }
		\hline
		Nombre & Tipo de dato & Valor Null & Primaria & For\'anea & Comentario \\ \hline
		UserIdUsuario & Cadena & Si & Si & No & \\ \hline
		RolFKIdRol & Entero & Si & No & Si & \\ \hline
		UserNombreUsuario & Cadena & Si & No & No & \\ \hline
		UserCorreo & Cadena & Si & No & No & \\ \hline
		UserNombres & Cadena & Si & No & No & \\ \hline
		UserApellidos & Cadena & Si & No & No & \\ \hline
	\end{tabular}
\end{table}

\begin{table}[ht]
	\caption{Tabla Semillero}
	\label{labelTableSemillero}
	\begin{tabular}{ |l|l|l|l|l|l| }
		\hline
		Nombre & Tipo de dato & Valor Null & Primaria & For\'anea & Comentario \\ \hline
		SmlrIdSemillero & Entero & Si & Si & No & \\ \hline 
		SmlrNombre & Cadena & Si & No & No & \\ \hline 
		SmlrDescripcion & Texto largo & Si & No & No & \\ \hline 
		SmlrRutaLogo & Texto largo & Si & No & No & \\ \hline 
		SmlrMision & Texto largo & Si & No & No & \\ \hline 
		SmlrVision & Texto largo & Si & No & No & \\ \hline 
		SmlrRuta & Cadena & Si & No & No & \\ \hline 
		SmlrSigla & Cadena & Si & No & No & \\ \hline 
		SmlrJustificacion & Texto largo & Si & No & No & \\ \hline 
		SmlrMetodologiaTrabajo & Texto largo & Si & No & No & \\ \hline 
		SmlrAreasTrabajo & Texto largo & Si & No & No & \\ \hline 	
	\end{tabular}
\end{table}

\begin{table}[ht]
	\caption{Tabla Integrantes Semillero}
	\label{labelTableIntegrantesSemillero}
	\begin{tabular}{ |l|l|l|l|l|l| }
		\hline
		Nombre & Tipo de dato & Valor Null & Primaria & For\'anea & Comentario \\ \hline
		IsmlIdIntegrante & Entero & Si & Si & No & \\ \hline 
		UserFKIdUsuario & Entero & Si & No & Si & \\ \hline 
		SmlrFKIdSemillero & Entero & Si & No & Si & \\ \hline 
		IsmlFecha & Fecha & Si & No & No & \\ \hline 
		IsmlEstado & Cadena & Si & No & No & \\ \hline	
	\end{tabular}
\end{table}

\begin{table}[ht]
	\caption{Tabla Rol}
	\label{labelTableRol}
	\begin{tabular}{ |l|l|l|l|l|l| }
		\hline
		Nombre & Tipo de dato & Valor Null & Primaria & For\'anea & Comentario \\ \hline
		RolIdRol & Entero & Si & Si & No & \\ \hline 
		RolNombreRol & Cadena & Si & No & No & \\ \hline	
	\end{tabular}
\end{table}

\begin{table}[ht]
	\caption{Tabla Permisos}
	\label{labelTablePermisos}
	\begin{tabular}{ |l|l|l|l|l|l| }
		\hline
		Nombre & Tipo de dato & Valor Null & Primaria & For\'anea & Comentario \\ \hline
		PrmsIdPermiso & Entero & Si & Si & No & \\ \hline 
		PrmsNombrePermiso & Cadena & Si & No & No & \\ \hline 
		PrmsIcono & Cadena & Si & No & No & \\ \hline 	
	\end{tabular}
\end{table}


\begin{table}[ht]
	\caption{Tabla Subpermisos}
	\label{labelTableSubpermisos}
	\begin{tabular}{ |l|l|l|l|l|l| }
		\hline
		Nombre & Tipo de dato & Valor Null & Primaria & For\'anea & Comentario \\ \hline
		SpmsIdSubpermiso & Entero & Si & Si & No & \\ \hline 
		SpmsNombreSubpermiso & Cadena & Si & No & No & \\ \hline 
		SpmsURL & Cadena & Si & No & No & \\ \hline 
		PrmsFKIdPermiso & Entero & Si & No & Si & \\ \hline 	
	\end{tabular}
\end{table}

\begin{table}[ht]
	\caption{Tabla Menú Usuario}
	\label{labelTableMenuUsuario}
	\begin{tabular}{ |l|l|l|l|l|l| }
		\hline
		Nombre & Tipo de dato & Valor Null & Primaria & For\'anea & Comentario \\ \hline
		RolFKIdRol & Entero & Si & Si & Si & \\ \hline 
		SpmsFKIdSubpermiso & Entero & Si & Si & Si & \\ \hline 
		MusrEstado & Cadena & Si & No & No & \\ \hline 	
	\end{tabular}
\end{table}

\begin{table}[ht]
	\caption{Tabla Repositorio}
	\label{labelTableRepositorio}
	\begin{tabular}{ |l|l|l|l|l|l| }
		\hline
		Nombre & Tipo de dato & Valor Null & Primaria & For\'anea & Comentario \\ \hline
		RepoIdRepositorio & Entero & Si & Si & No & \\ \hline 
		RepoNombre & Cadena & Si & No & No & \\ \hline 
		RepoDescripcion & Texto largo & Si & No & No & \\ \hline 
		RepoRuta & Texto largo & Si & No & No & \\ \hline 
		RepoEstado & Cadena & Si & No & No & \\ \hline 	
	\end{tabular}
\end{table}


\begin{table}[ht]
	\caption{Tabla Documento}
	\label{labelTableDocumento}
	\begin{tabular}{ |l|l|l|l|l|l| }
		\hline
		Nombre & Tipo de dato & Valor Null & Primaria & For\'anea & Comentario \\ \hline
		DocuIdDocumento & Entero & Si & Si & No & \\ \hline 
		DocuNombre & Cadena & Si & No & No & \\ \hline 
		DocuRuta & Texto largo & Si & No & No & \\ \hline 
		DocuFecha & Fecha & Si & No & No & \\ \hline 
		DocuEstado & Cadena & No & No & No & \\ \hline 
		UserFKIdUsuario & BI Entero & Si & No & Si & \\ \hline 
		RepoFKIdRepositorio & Entero & Si & No & Si & \\ \hline 	
	\end{tabular}
\end{table}

\begin{table}[ht]
	\caption{Tabla Correos Enviados}
	\label{labelTableCorreosEnviados}
	\begin{tabular}{ |l|l|l|l|l|l| }
		\hline
		Nombre & Tipo de dato & Valor Null & Primaria & For\'anea & Comentario \\ \hline
		CrnvIdCorreoEnviado & Entero & Si & Si & No & \\ \hline 
		CrnvTipo & Cadena & Si & No & No & \\ \hline 
		CrnvDescripcion & Texto largo & Si & No & No & \\ \hline 
		CrnvEstadoEnvio & VARCHAR(50) & Si & No & No & \\ \hline 
		CrnvFechaEnvio & Fecha & Si & No & No & \\ \hline 
		SmlrFKIdSemillero & Entero & Si & No & Si & \\ \hline 
		UserFKDestinatario & Cadena & Si & No & Si & \\ \hline 	
	\end{tabular}
\end{table}


\begin{table}[ht]
	\caption{Tabla Acceso Sistema}
	\label{labelTableAccesoSistema}
	\begin{tabular}{ |l|l|l|l|l|l| }
		\hline
		Nombre & Tipo de dato & Valor Null & Primaria & For\'anea & Comentario \\ \hline
		AstmIdAcceso & BI Entero & Si & Si & No & \\ \hline 
		AstmFecha & Fecha & Si & No & No & \\ \hline 
		AstmIP & Cadena & Si & No & No & \\ \hline 
		UserFKIdUsuario & BI Entero & Si & No & Si & \\ \hline 	
	\end{tabular}
\end{table}


\begin{table}[ht]
	\caption{Tabla Solicitudes}
	\label{labelTableSolicitudes}
	\begin{tabular}{ |l|l|l|l|l|l| }
		\hline
		Nombre & Tipo de dato & Valor Null & Primaria & For\'anea & Comentario \\ \hline
		SlctIdSolicitud & Entero & Si & Si & No & \\ \hline 
		CursIdCurso & Entero & Si & No & No & \\ \hline 
		UserFKEstudiante & BI Entero & Si & No & Si & \\ \hline 
		UserFKAdmin & BI Entero & Si & No & Si & \\ \hline 
		SlctFechaSolicitud & Fecha & Si & No & No & \\ \hline 
		SlctFechaRespuesta & Fecha & No & No & No & \\ \hline 
		SlctEstado & Cadena & No & No & No & \\ \hline 
		SlctComentario & Texto largo & Si & No & No & \\ \hline 	
	\end{tabular}
\end{table}


\begin{table}[ht]
	\caption{Tabla Grupo Investigaci\'on}
	\label{labelTableGrupoInvestigacion}
	\begin{tabular}{ |l|l|l|l|l|l| }
		\hline
		Nombre & Tipo de dato & Valor Null & Primaria & For\'anea & Comentario \\ \hline
		GrivIdGrupoInv & Entero & Si & Si & No & \\ \hline 
		GrivFechaCreacion & Fecha & Si & No & No & \\ \hline 
		GrivNombre & Cadena & Si & No & No & \\ \hline 
		GrivMision & Texto largo & Si & No & No & \\ \hline 
		GrivVision & Texto largo & Si & No & No & \\ \hline 	
	\end{tabular}
\end{table}


\begin{table}[ht]
	\caption{Tabla Programa}
	\label{labelTablePrograma}
	\begin{tabular}{ |l|l|l|l|l|l| }
		\hline
		Nombre & Tipo de dato & Valor Null & Primaria & For\'anea & Comentario \\ \hline
		ProgIdPrograma & Entero & Si & Si & No & \\ \hline 
		ProgNombre & Cadena & Si & No & No & \\ \hline 
		ProgCodigo & Cadena & No & No & No & \\ \hline 
		FacuFKIdFacultad & Entero & Si & No & Si & \\ \hline 	
	\end{tabular}
\end{table}


\begin{table}[ht]
	\caption{Tabla Programas Semillero}
	\label{labelTableProgramasSemillero}
	\begin{tabular}{ |l|l|l|l|l|l| }
		\hline
		Nombre & Tipo de dato & Valor Null & Primaria & For\'anea & Comentario \\ \hline
		ProgFKIdPrograma & Entero & Si & Si & Si & \\ \hline 
		SmlrFKIdSemillero & Entero & Si & Si & Si & \\ \hline 
		PgsmEstado & ENUM('T', 'F') & Si & No & No & \\ \hline 	
	\end{tabular}
\end{table}


\begin{table}[ht]
	\caption{Tabla Programas Grupos Inv}
	\label{labelTableProgramasGruposInv}
	\begin{tabular}{ |l|l|l|l|l|l| }
		\hline
		Nombre & Tipo de dato & Valor Null & Primaria & For\'anea & Comentario \\ \hline
		GrivFKIdGrupoInv & Entero & Si & Si & Si & \\ \hline 
		ProgFKIdPrograma & Entero & Si & Si & Si & \\ \hline 
		PggrEstado & ENUM('T', 'F') & Si & No & No & \\ \hline 
		PggrFecha & Fecha & Si & No & No & \\ \hline 
		PggrCartaAval & Texto largo & Si & No & No & \\ \hline 	
	\end{tabular}
\end{table}


\begin{table}[ht]
	\caption{Tabla Proyectos}
	\label{labelTableProyectos}
	\begin{tabular}{ |l|l|l|l|l|l| }
		\hline
		Nombre & Tipo de dato & Valor Null & Primaria & For\'anea & Comentario \\ \hline
		ProyIdProyecto & Entero & Si & Si & No & \\ \hline 
		ProyNombre & Texto largo & Si & No & No & \\ \hline 
		ProyDescripcion & Texto largo & Si & No & No & \\ \hline 
		ProyTipo & Cadena & Si & No & No & \\ \hline 
		ProSitado & Cadena & Si & No & No & \\ \hline 
		LinvFKIdLineaInv & Entero & Si & No & Si & \\ \hline 
		ConvFKIdConvocatoria & Entero & Si & No & Si & \\ \hline 
		GrivFKIdGrupoInv & Entero & Si & No & Si & \\ \hline 
		ProgFKIdPrograma & Entero & Si & No & Si & \\ \hline 
		ProyFormulacionProblema & Texto largo & Si & No & No & \\ \hline 
		ProyJustificacion & Texto largo & Si & No & No & \\ \hline 
		ProyMarcoTeorico & Texto largo & Si & No & No & \\ \hline 
		ProyMetodologia & Texto largo & Si & No & No & \\ \hline 
		ProyBibliografia & Texto largo & Si & No & No & \\ \hline 
		ProyAnexos & Texto largo & Si & No & No & \\ \hline 	
	\end{tabular}
\end{table}


\begin{table}[ht]
	\caption{Tabla Rol Semillero}
	\label{labelTableRolSemillero}
	\begin{tabular}{ |l|l|l|l|l|l| }
		\hline
		Nombre & Tipo de dato & Valor Null & Primaria & For\'anea & Comentario \\ \hline
		RsmlIdRolSemillero & Entero & Si & Si & No & \\ \hline 
		RsmlNombre & Cadena & Si & No & No & \\ \hline 
		RsmlEstado & Cadena & Si & No & No & \\ \hline 
		SmlrFKIdSemillero & Entero & Si & No & Si & \\ \hline 	
	\end{tabular}
\end{table}


\begin{table}[ht]
	\caption{Tabla Par Evaluador}
	\label{labelTableParEvaluador}
	\begin{tabular}{ |l|l|l|l|l|l| }
		\hline
		Nombre & Tipo de dato & Valor Null & Primaria & For\'anea & Comentario \\ \hline
		PaevIdParEvaluador & Entero & Si & Si & No & \\ \hline 
		PaevFecha & Fecha & Si & No & No & \\ \hline 
		PaevEstado & Cadena & Si & No & No & \\ \hline 
		PaevTipo & Cadena & Si & No & No & \\ \hline 
		UserFKParEvaluador & Cadena & Si & No & Si & \\ \hline 	
	\end{tabular}
\end{table}


\begin{table}[ht]
	\caption{Tabla Actividad Semillero}
	\label{labelTableActividadSemillero}
	\begin{tabular}{ |l|l|l|l|l|l| }
		\hline
		Nombre & Tipo de dato & Valor Null & Primaria & For\'anea & Comentario \\ \hline
		ActsIdActividad & Entero & Si & Si & No & \\ \hline 
		ActsDescripcion & Texto largo & Si & No & No & \\ \hline 
		ActsFechaCreacion & Fecha & Si & No & No & \\ \hline 
		ActsEstado & Cadena & Si & No & No & \\ \hline 
		ActsFechaInicio & DATE & Si & No & No & \\ \hline 
		ActsFechaFin & DATE & Si & No & No & \\ \hline 
		MetsFKIdMetaSemillero & Entero & Si & No & Si & \\ \hline 	
	\end{tabular}
\end{table}


\begin{table}[ht]
	\caption{Tabla Convocatoria}
	\label{labelTableConvocatoria}
	\begin{tabular}{ |l|l|l|l|l|l| }
		\hline
		Nombre & Tipo de dato & Valor Null & Primaria & For\'anea & Comentario \\ \hline
		ConvIdConvocatoria & Entero & Si & Si & No & \\ \hline 
		ConvNombre & Texto largo & Si & No & No & \\ \hline 
		ConvDescripcion & Texto largo & Si & No & No & \\ \hline 
		ConvObservaciones & Texto largo & No & No & No & \\ \hline 
		ConvFechaCreacion & Fecha & Si & No & No & \\ \hline 
		ConvFechaInicio & DATE & Si & No & No & \\ \hline 
		ConvFechaFin & DATE & Si & No & No & \\ \hline 
		ConvEstado & Cadena & Si & No & No & \\ \hline 
		TconFKIdTipoConvocatoria & Entero & Si & No & Si & \\ \hline 
		ConvOferente & Cadena & Si & No & No & \\ \hline 
		ConvDestinatario & Cadena & Si & No & No & \\ \hline 	
	\end{tabular}
\end{table}


\begin{table}[ht]
	\caption{Tabla Banco Preliminar}
	\label{labelTableBancopreliminar}
	\begin{tabular}{ |l|l|l|l|l|l| }
		\hline
		Nombre & Tipo de dato & Valor Null & Primaria & For\'anea & Comentario \\ \hline
		BpreIdBancoPreliminar & Entero & Si & Si & No & \\ \hline 
		BpreDescripcion & Cadena & Si & No & No & \\ \hline 
		BpreEstado & Cadena & Si & No & No & \\ \hline 
		ConvFKIdConvocatoria & Entero & Si & No & Si & \\ \hline 
		BpreFecha & Fecha & Si & No & No & \\ \hline 	
	\end{tabular}
\end{table}


\begin{table}[ht]
	\caption{Tabla Proyectos preliminares}
	\label{labelTableProyectospreliminares}
	\begin{tabular}{ |l|l|l|l|l|l| }
		\hline
		Nombre & Tipo de dato & Valor Null & Primaria & For\'anea & Comentario \\ \hline
		PpreIdProyectoPreliminar & Entero & Si & Si & No & \\ \hline 
		ProyFKIdProyecto & Entero & Si & No & Si & \\ \hline 
		BpreFKIdBancoPreliminar & Entero & Si & No & Si & \\ \hline 
		PpreEstado & Cadena & Si & No & No & \\ \hline 
		PpreFecha & Fecha & Si & No & No & \\ \hline 	
	\end{tabular}
\end{table}


\begin{table}[ht]
	\caption{Tabla Banco definitivo}
	\label{labelTableBancodefinitivo}
	\begin{tabular}{ |l|l|l|l|l|l| }
		\hline
		Nombre & Tipo de dato & Valor Null & Primaria & For\'anea & Comentario \\ \hline
		BdefIdBancoDefinitivo & Entero & Si & Si & No & \\ \hline 
		BdefEstado & Cadena & Si & No & No & \\ \hline 
		BdefFecha & Fecha & Si & No & No & \\ \hline 
		ConvFKIdConvocatoria & Entero & Si & No & Si & \\ \hline 	
	\end{tabular}
\end{table}


\begin{table}[ht]
	\caption{Tabla Proyectos definitivos}
	\label{labelTableProyectosdefinitivos}
	\begin{tabular}{ |l|l|l|l|l|l| }
		\hline
		Nombre & Tipo de dato & Valor Null & Primaria & For\'anea & Comentario \\ \hline
		PdefIdProyectoDefinitivo & Entero & Si & Si & No & \\ \hline 
		PdefEstado & Cadena & Si & No & No & \\ \hline 
		PdefFecha & Fecha & Si & No & No & \\ \hline 
		BdefFKIdBancoDefinitivo & Entero & Si & No & Si & \\ \hline 
		ProyFKIdProyecto & Entero & Si & No & Si & \\ \hline 	
	\end{tabular}
\end{table}


\begin{table}[ht]
	\caption{Tabla Evaluaci\'on}
	\label{labelTableEvaluacion}
	\begin{tabular}{ |l|l|l|l|l|l| }
		\hline
		Nombre & Tipo de dato & Valor Null & Primaria & For\'anea & Comentario \\ \hline
		EvlcIdEvaluacion & BI Entero & Si & Si & No & \\ \hline 
		EvlcDescripcion & Cadena & Si & No & No & \\ \hline 
		EvlcFecha & Fecha & Si & No & No & \\ \hline 
		EvlcEstado & Cadena & Si & No & No & \\ \hline 
		EvlcResultado & Cadena & Si & No & No & \\ \hline 
		AparFKIdAsignacion & Entero & Si & No & Si & \\ \hline 	
	\end{tabular}
\end{table}


\begin{table}[ht]
	\caption{Tabla Tipo Convocatoria}
	\label{labelTableTipoConvocatoria}
	\begin{tabular}{ |l|l|l|l|l|l| }
		\hline
		Nombre & Tipo de dato & Valor Null & Primaria & For\'anea & Comentario \\ \hline
		TconIdTipoConvocatoria & Entero & Si & Si & No & \\ \hline 
		TconNombre & Cadena & Si & No & No & \\ \hline 
		TconEstado & Cadena & Si & No & No & \\ \hline 	
	\end{tabular}
\end{table}


\begin{table}[ht]
	\caption{Tabla Linea Investigaci\'on}
	\label{labelTableLineaInvestigacion}
	\begin{tabular}{ |l|l|l|l|l|l| }
		\hline
		Nombre & Tipo de dato & Valor Null & Primaria & For\'anea & Comentario \\ \hline
		LinvIdLineaInv & Entero & Si & Si & No & \\ \hline 
		LinvNombreLinea & Cadena & Si & No & No & \\ \hline 
		LinvDescripcion & Texto largo & Si & No & No & \\ \hline 
		ProgFKIdPrograma & Entero & Si & No & Si & \\ \hline 	
	\end{tabular}
\end{table}


\begin{table}[ht]
	\caption{Tabla Facultad}
	\label{labelTableFacultad}
	\begin{tabular}{ |l|l|l|l|l|l| }
		\hline
		Nombre & Tipo de dato & Valor Null & Primaria & For\'anea & Comentario \\ \hline
		FacuIdFacultad & Entero & Si & Si & No & \\ \hline 
		FacuNombre & Cadena & Si & No & No & \\ \hline 
		FacuEstado & Cadena & Si & No & No & \\ \hline 	
	\end{tabular}
\end{table}


\begin{table}[ht]
	\caption{Tabla Evaluador Interno}
	\label{labelTableEvaluadorInterno}
	\begin{tabular}{ |l|l|l|l|l|l| }
		\hline
		Nombre & Tipo de dato & Valor Null & Primaria & For\'anea & Comentario \\ \hline
		EvaiIdEvaluadorInterno & Entero & Si & Si & No & \\ \hline 
		ProgFKIdPrograma & Entero & Si & No & Si & \\ \hline 
		PaevFKIdParEvaluador & Entero & Si & No & Si & \\ \hline 	
	\end{tabular}
\end{table}


\begin{table}[ht]
	\caption{Tabla Asignaci\'on Par Evaluaci\'on}
	\label{labelTableAsignacionParEvaluacion}
	\begin{tabular}{ |l|l|l|l|l|l| }
		\hline
		Nombre & Tipo de dato & Valor Null & Primaria & For\'anea & Comentario \\ \hline
		AparIdAsignacion & Entero & Si & Si & No & \\ \hline 
		ProyFKIdProyecto & Entero & Si & No & Si & \\ \hline 
		PaevFKIdParEvaluador & Entero & Si & No & Si & \\ \hline 
		AparEstado & Cadena & Si & No & No & \\ \hline 
		AparFecha & Fecha & Si & No & No & \\ \hline 
		UserFKAsignadoPor & Cadena & Si & No & Si & \\ \hline 	
	\end{tabular}
\end{table}


\begin{table}[ht]
	\caption{Tabla Preevaluacion}
	\label{labelTablePreevaluacion}
	\begin{tabular}{ |l|l|l|l|l|l| }
		\hline
		Nombre & Tipo de dato & Valor Null & Primaria & For\'anea & Comentario \\ \hline
		PrevIdPreevaluacion & Entero & Si & Si & No & \\ \hline 
		PrevDescripcion & Texto largo & Si & No & No & \\ \hline 
		PrevFecha & Fecha & Si & No & No & \\ \hline 
		PrevEstado & Cadena & Si & No & No & \\ \hline 
		PrevResultado & Cadena & Si & No & No & \\ \hline 
		InstFKIdInstancia & Entero & Si & No & Si & \\ \hline 
		ProyFKIdProyecto & Entero & Si & No & Si & \\ \hline 
		UserFKIdUsuario & Cadena & Si & No & Si & \\ \hline 	
	\end{tabular}
\end{table}


\begin{table}[ht]
	\caption{Tabla Instancia}
	\label{labelTableInstancia}
	\begin{tabular}{ |l|l|l|l|l|l| }
		\hline
		Nombre & Tipo de dato & Valor Null & Primaria & For\'anea & Comentario \\ \hline
		InstIdInstancia & Entero & Si & Si & No & \\ \hline 
		InstNombre & Cadena & Si & No & No & \\ \hline 
		InstDescripcion & Texto largo & Si & No & No & \\ \hline 
		InstEstado & Cadena & Si & No & No & \\ \hline 	
	\end{tabular}
\end{table}


\begin{table}[ht]
	\caption{Tabla Integrantes Grupo}
	\label{labelTableIntegrantesGrupo}
	\begin{tabular}{ |l|l|l|l|l|l| }
		\hline
		Nombre & Tipo de dato & Valor Null & Primaria & For\'anea & Comentario \\ \hline
		IgruIdIntegrante & Entero & Si & Si & No & \\ \hline 
		UserFKIdUsuario & Cadena & Si & No & Si & \\ \hline 
		GrivFKIdGrupoInv & Entero & Si & No & Si & \\ \hline 
		IgruEstado & Cadena & Si & No & No & \\ \hline 
		IgruFecha & Fecha & Si & No & No & \\ \hline 	
	\end{tabular}
\end{table}


\begin{table}[ht]
	\caption{Tabla Proyectos Semilleros}
	\label{labelTableProyectosSemilleros}
	\begin{tabular}{ |l|l|l|l|l|l| }
		\hline
		Nombre & Tipo de dato & Valor Null & Primaria & For\'anea & Comentario \\ \hline
		PrsmIdProyecto & Entero & Si & Si & No & \\ \hline 
		PrsmFecha & Fecha & Si & No & No & \\ \hline 
		PrsmEstado & Cadena & Si & No & No & \\ \hline 
		ProyFKIdProyecto & Entero & Si & No & Si & \\ \hline 
		SmlrFKIdSemillero & Entero & Si & No & Si & \\ \hline 	
	\end{tabular}
\end{table}


\begin{table}[ht]
	\caption{Tabla Proyectos Grupos}
	\label{labelTableProyectosGrupos}
	\begin{tabular}{ |l|l|l|l|l|l| }
		\hline
		Nombre & Tipo de dato & Valor Null & Primaria & For\'anea & Comentario \\ \hline
		PrgrIdProyecto & Entero & Si & Si & No & \\ \hline 
		PrgrFecha & Fecha & Si & No & No & \\ \hline 
		PrgrEstado & Cadena & Si & No & No & \\ \hline 
		ProyFKIdProyecto & Entero & Si & No & Si & \\ \hline 
		GrivFKIdGrupoInv & Entero & Si & No & Si & \\ \hline 
		PrgrLugarEjecucion & Cadena & Si & No & No & \\ \hline 
		MpioFKIdMpio & Entero & Si & No & Si & \\ \hline 
		PrgrDuracion & DOUBLE & Si & No & No & \\ \hline 
		TinvFKIdTipo & Entero & Si & No & Si & \\ \hline 
		PrgrValorTotal & BI Entero & Si & No & No & \\ \hline 
		PrgrValorSolicitado & BI Entero & Si & No & No & \\ \hline 
		PrgrValorContrapartida & BI Entero & Si & No & No & \\ \hline 
		PrgrTotalInvestigadores & Entero & Si & No & No & \\ \hline 
		PrgrTipo & Cadena & Si & No & No & \\ \hline 
		PrgrResultadosEsperados & Texto largo & Si & No & No & \\ \hline 
		PrgrEstraComunicacion & Texto largo & Si & No & No & \\ \hline 
		PrgrPalabrasClave & Texto largo & Si & No & No & \\ \hline 
		& Cadena & No & No & No & \\ \hline 	
	\end{tabular}
\end{table}


\begin{table}[ht]
	\caption{Tabla Roles Integrantes Semillero}
	\label{labelTableRolesIntegrantesSemillero}
	\begin{tabular}{ |l|l|l|l|l|l| }
		\hline
		Nombre & Tipo de dato & Valor Null & Primaria & For\'anea & Comentario \\ \hline
		RismIdRol & Entero & Si & Si & No & \\ \hline 
		RismFecha & Fecha & Si & No & No & \\ \hline 
		RismEstado & Cadena & Si & No & No & \\ \hline 
		IsmlFKIdIntegrante & Entero & Si & No & Si & \\ \hline 
		RsmlFKIdRolSemillero & Entero & Si & No & Si & \\ \hline 	
	\end{tabular}
\end{table}


\begin{table}[ht]
	\caption{Tabla Rol Grupo}
	\label{labelTableRolGrupo}
	\begin{tabular}{ |l|l|l|l|l|l| }
		\hline
		Nombre & Tipo de dato & Valor Null & Primaria & For\'anea & Comentario \\ \hline
		RgruIdRolGrupo & Entero & Si & Si & No & \\ \hline 
		RgruNombre & Cadena & Si & No & No & \\ \hline 
		RgruEstado & Cadena & Si & No & No & \\ \hline 
		GrivFKIdGrupoInv & Entero & Si & No & Si & \\ \hline 	
	\end{tabular}
\end{table}


\begin{table}[ht]
	\caption{Tabla Roles Integrantes Grupo}
	\label{labelTableRolesIntegrantesGrupo}
	\begin{tabular}{ |l|l|l|l|l|l| }
		\hline
		Nombre & Tipo de dato & Valor Null & Primaria & For\'anea & Comentario \\ \hline
		RigrIdRol & Entero & Si & Si & No & \\ \hline 
		RigrFecha & Fecha & Si & No & No & \\ \hline 
		RigrEstado & Cadena & Si & No & No & \\ \hline 
		IgruFKIdIntegrante & Entero & Si & No & Si  & \\ \hline 
		RgruFKIdRolGrupo & Entero & Si & No & Si & \\ \hline 	
	\end{tabular}
\end{table}


\begin{table}[ht]
	\caption{Tabla Objetivos Proyecto Semillero}
	\label{labelTableObjetivosProyectoSemillero}
	\begin{tabular}{ |l|l|l|l|l|l| }
		\hline
		Nombre & Tipo de dato & Valor Null & Primaria & For\'anea & Comentario \\ \hline
		ObpsIdObjetivo & Entero & Si & Si & No & \\ \hline 
		ObpsID & Cadena & Si & No & No & \\ \hline 
		ObpsDescripcion & Cadena & Si & No & No & \\ \hline 
		ObpsFecha & Fecha & Si & No & No & \\ \hline 
		ObpsEstado & Cadena & Si & No & No & \\ \hline 
		ObpsTipo & Cadena & Si & No & No & \\ \hline 
		PrsmFKIdProyecto & Entero & Si & No & Si & \\ \hline 	
	\end{tabular}
\end{table}


\begin{table}[ht]
	\caption{Tabla Informe Parcial}
	\label{labelTableInformeParcial}
	\begin{tabular}{ |l|l|l|l|l|l| }
		\hline
		Nombre & Tipo de dato & Valor Null & Primaria & For\'anea & Comentario \\ \hline
		InpaIdInformeParcial & Entero & Si & Si & No & \\ \hline 
		InpaDescripcion & Texto largo & Si & No & No & \\ \hline 
		InpaFecha & Fecha & Si & No & No & \\ \hline 
		InpaEstado & Cadena & Si & No & No & \\ \hline 
		InpaMetodologia & Texto largo & Si & No & No & \\ \hline 
		InpaResultados & Texto largo & Si & No & No & \\ \hline 
		InpaImpacto & Texto largo & Si & No & No & \\ \hline 
		InpaEstraDivulgacion & Texto largo & Si & No & No & \\ \hline 
		ProyFKIdProyecto & Entero & Si & No & Si & \\ \hline 
		InpaDocSoporte & Texto largo & Si & No & No & \\ \hline 
		InpaPresupuesto & Texto largo & Si & No & No & \\ \hline 	
	\end{tabular}
\end{table}


\begin{table}[ht]
	\caption{Tabla Integrantes Proyecto Semillero}
	\label{labelTableIntegrantesProyectoSemillero}
	\begin{tabular}{ |l|l|l|l|l|l| }
		\hline
		Nombre & Tipo de dato & Valor Null & Primaria & For\'anea & Comentario \\ \hline
		IprsIdIntegrante & Entero & Si & Si & No & \\ \hline 
		IprsFecha & Fecha & Si & No & No & \\ \hline 
		IprsEstado & Cadena & Si & No & No & \\ \hline 
		RolpFKIdRolProyecto & Entero & Si & No & Si & \\ \hline 
		IsmlFKIdIntegrante & Entero & Si & No & Si & \\ \hline 
		PrsmFKIdProyecto & Entero & Si & No & Si &  \\ \hline 	
	\end{tabular}
\end{table}


\begin{table}[ht]
	\caption{Tabla Rol Proyecto}
	\label{labelTableRolProyecto}
	\begin{tabular}{ |l|l|l|l|l|l| }
		\hline
		Nombre & Tipo de dato & Valor Null & Primaria & For\'anea & Comentario \\ \hline
		RolpIdRolProyecto & Entero & Si & Si & No & \\ \hline 
		RolpNombre & Cadena & Si & No & No & \\ \hline 
		RolpEstado & Cadena & Si & No & No & \\ \hline 	
	\end{tabular}
\end{table}


\begin{table}[ht]
	\caption{Tabla Informe Final}
	\label{labelTableInformeFinal}
	\begin{tabular}{ |l|l|l|l|l|l| }
		\hline
		Nombre & Tipo de dato & Valor Null & Primaria & For\'anea & Comentario \\ \hline
		InfiIdInformeFinal & Entero & Si & Si & No & \\ \hline 
		InfiProblemaPlanteado & Texto largo & Si & No & No & \\ \hline 
		InfiHipotesis & Texto largo & Si & No & No & \\ \hline 
		InfiAvanceObtenido & Texto largo & Si & No & No & \\ \hline 
		InfiRetosPlanteados & Texto largo & Si & No & No & \\ \hline 
		InfiGradoComprobacion & Texto largo & Si & No & No & \\ \hline 
		InfiGestionProyecto & Texto largo & Si & No & No & \\ \hline 
		InfiObservaciones & Texto largo & Si & No & No & \\ \hline 
		InfiRecomendaciones & Texto largo & Si & No & No & \\ \hline 
		InfiActores & Texto largo & Si & No & No & \\ \hline 	
	\end{tabular}
\end{table}


\begin{table}[ht]
	\caption{Tabla Informe Final}
	\label{labelTableInformeFinal}
	\begin{tabular}{ |l|l|l|l|l|l| }
		\hline
		Nombre & Tipo de dato & Valor Null & Primaria & For\'anea & Comentario \\ \hline
		InfiIdInformeFinal & Entero & Si & Si & No & \\ \hline 
		InfiProblemaPlanteado & Texto largo & Si & No & No & \\ \hline 
		InfiHipotesis & Texto largo & Si & No & No & \\ \hline 
		InfiAvanceObtenido & Texto largo & Si & No & No & \\ \hline 
		InfiRetosPlanteados & Texto largo & Si & No & No & \\ \hline 
		InfiGradoComprobacion & Texto largo & Si & No & No & \\ \hline 
		InfiGestionProyecto & Texto largo & Si & No & No & \\ \hline 
		InfiObservaciones & Texto largo & Si & No & No & \\ \hline 
		InfiRecomendaciones & Texto largo & Si & No & No & \\ \hline 
		InfiActores & Texto largo & Si & No & No & \\ \hline 	
	\end{tabular}
\end{table}


\begin{table}[ht]
	\caption{Tabla Cronograma Convocatoria}
	\label{labelTableCronogramaConvocatoria}
	\begin{tabular}{ |l|l|l|l|l|l| }
		\hline
		Nombre & Tipo de dato & Valor Null & Primaria & For\'anea & Comentario \\ \hline
		CrcvIdCronograma & Entero & Si & Si & No & \\ \hline 
		CrcvDescripcion & Texto largo & Si & No & No & \\ \hline 
		CrcvFechaCreacion & Fecha & Si & No & No & \\ \hline 
		CrcvFechaActividad & Fecha & Si & No & No & \\ \hline 
		CrcvEstado & Cadena & No & No & No & \\ \hline 
		ConvFKIdConvocatoria & Entero & Si & No & Si & \\ \hline 	
	\end{tabular}
\end{table}

\begin{table}[ht]
	\caption{Tabla Soportes Cronograma}
	\label{labelTableSoportesCronograma}
	\begin{tabular}{ |l|l|l|l|l|l| }
		\hline
		Nombre & Tipo de dato & Valor Null & Primaria & For\'anea & Comentario \\ \hline
		SocrIdSoporte & Entero & Si & Si & No & \\ \hline 
		SocrDescripcion & Texto largo & Si & No & No & \\ \hline 
		SocrRutaAdjunto & Cadena & Si & No & No & \\ \hline 
		SocrEstado & Cadena & Si & No & No & \\ \hline 
		SocrFecha & Fecha & Si & No & No & \\ \hline 
		CrcvFKIdCronograma & Entero & Si & No & Si &  \\ \hline	
	\end{tabular}
\end{table}

\begin{table}[ht]
	\caption{Tabla Producci\'on Grupo}
	\label{labelTableProduccionGrupo}
	\begin{tabular}{ |l|l|l|l|l|l| }
		\hline
		Nombre & Tipo de dato & Valor Null & Primaria & For\'anea & Comentario \\ \hline
		PgruIdProduccion & Entero & Si & Si & No & \\ \hline 
		Pgru & Cadena & No & No & No & \\ \hline 	
	\end{tabular}
\end{table}

\begin{table}[ht]
	\caption{Tabla Unidades Acad\'emicas}
	\label{labelTableUnidadesAcademicas}
	\begin{tabular}{ |l|l|l|l|l|l| }
		\hline
		Nombre & Tipo de dato & Valor Null & Primaria & For\'anea & Comentario \\ \hline		
		UnacIdUnidad & Entero & Si & Si & No & \\ \hline 
		UnacNombre & Cadena & Si & No & No & \\ \hline 
		UnacEstado & Cadena & Si & No & No & \\ \hline 
		UnacFechaCreacion & Fecha & Si & No & No & \\ \hline	
	\end{tabular}
\end{table}
	\chapter{Licitaci\'on de requisitos}

\section{Objetivos del sistem\'a}

\begin{table}[H]
	\centering
	\caption{Requisito - Gestionar convocatoria}
	\label{label_req_gconv}
	\begin{tabular}{|l|p{13cm}|}
		\hline
		OBJ-001 & \begin{tabular}[c]{@{}l@{}}
			Gestionar convocatorias
		\end{tabular} \\ \hline
		Versi\'on & 002 (2016-08-03) \\ \hline
		Autores & \begin{tabular}[c]{@{}l@{}}
			Angie Zuleta Cardona (Universidad de la Amazonia) \\
			Derly Viviana  Murcia Serrano (Universidad de la Amazonia)\\
			Mateo Ceballos Bermudez (Universidad de la Amazonia) \\
			Nicol Dayana Endo Ruiz (Universidad de la Amazonia)
		\end{tabular} \\ \hline
		Fuentes & \begin{tabular}[c]{@{}l@{}}
			Alberto Fajardo Oliveros (Uniamazonia) \\
			Dora Lida Suarez Castro (Uniamazonia)
		\end{tabular} \\ \hline
		Descripci\'on & \begin{tabular}[c]{@{}l@{}}
			El sistema deber\'a gestionar la informaci\'on \\ 
			respecto a las distintas convocatorias en el \\ 
			marco de semilleros y grupos de investigaci\'on.
		\end{tabular} \\ \hline
		SubObjetivos & \\ \hline
		Importancia & Vital \\ \hline
		Urgencia & Inmediatamente \\ \hline
		Estado & En Construcci\'on \\ \hline
		Estabilidad & Alta \\ \hline
		Comentarios & Ninguno \\ \hline
	\end{tabular}
\end{table}


\begin{table}[H]
	\centering
	\caption{Requisito - Gestionar propuestas de proyectos}
	\label{label_req_gpdp}
	\begin{tabular}{|l|l|}
		\hline
		OBJ-002      & \begin{tabular}[c]{@{}l@{}}
			Gestionar propuestas de proyectos
		\end{tabular} \\ \hline
		Versi\'on    & 003 (2016-08-03) \\ \hline
		Autores      & \begin{tabular}[c]{@{}l@{}}
			Angie Zuleta Cardona (Universidad de la Amazonia)\\
			Derly Viviana  Murcia Serrano (Universidad de la Amazonia) \\ 
			Mateo Ceballos Bermudez (Universidad de la Amazonia) \\ 
			Nicol Dayana Endo Ruiz (Universidad de la Amazonia)
		\end{tabular} \\ \hline
		Fuentes      & \begin{tabular}[c]{@{}l@{}}
			Alberto Fajardo Oliveros (Uniamazonia) \\
			Dora Lida Suarez Castro (Uniamazonia)
		\end{tabular} \\ \hline
		Descripci\'on  & \begin{tabular}[c]{@{}l@{}}
			El sistema debe gestionar la informaci\'on que\\
			respecta a cada una de las propuestas que se \\
			presentan, sin importar el origen de estas.
		\end{tabular} \\  \hline
		SubObjetivos & \\  \hline
		Importancia  & Vital  \\ \hline
		Urgencia     & Inmediatamente \\
		Estado       & En Construcci\'on \\
		Estabilidad  & Alta \\ \hline
		Comentarios  & Ninguno \\  \hline
	\end{tabular}
\end{table}

\begin{table}[H]
	\centering
	\caption{Requisito - Gesti\'on de informaci\'on proyectos de investigaci\'on}
	\label{label_req_gipi}
	\begin{tabular}{|l|l|}
		\hline
		Autores      & \begin{tabular}[c]{@{}l@{}}
			Angie Zuleta Cardona (Universidad de la Amazonia)\\
			Derly Viviana  Murcia Serrano (Universidad de la Amazonia)\\
			Mateo Ceballos Bermudez (Universidad de la Amazonia)\\
			Nicol Dayana Endo Ruiz (Universidad de la Amazonia)
		\end{tabular} \\ \hline
		Fuentes      & \begin{tabular}[c]{@{}l@{}}
			Alberto Fajardo Oliveros (Uniamazonia) \\
			Dora Lida Suarez Castro (Uniamazonia)
		\end{tabular} \\ \hline
		Descripci\'on  & \begin{tabular}[c]{@{}l@{}}
			El sistema permitir\'a gestionar la informaci\'on\\
			correspondiente a cada proyecto de investigaci\'on
		\end{tabular} \\ \hline
		SubObjetivos & \\ \hline
		Importancia  & Vital \\ \hline
		Urgencia     & Inmediatamente \\ \hline
		Estado       & En Construcci\'on \\ \hline
		Estabilidad  & Alta \\ \hline
		Comentarios & \\ \hline
	\end{tabular}
\end{table}


\begin{table}[H]
	\centering
	\caption{Requisito - Gestionar informaci\'on de Grupos de Investigaci\'on}
	\label{label_req_gigi}
	\begin{tabular}{|l|l|}
		\hline
		OBJ-004      & \begin{tabular}[c]{@{}l@{}}
			Gestionar informaci\'on de Grupos de Investigaci\'on
		\end{tabular} \\ \hline
		Versi\'on    & 001 (2016-08-03) \\ \hline
		Autores      & \begin{tabular}[c]{@{}l@{}}
			Angie Zuleta Cardona (Universidad de la Amazonia)\\   
			Derly Viviana  Murcia Serrano (Universidad de la Amazonia)\\
			Mateo Ceballos Bermudez (Universidad de la Amazonia)\\
			Nicol Dayana Endo Ruiz (Universidad de la Amazonia)
		\end{tabular} \\ \hline
		Fuentes      & \begin{tabular}[c]{@{}l@{}}
			Alberto Fajardo Oliveros (Uniamazonia)\\
			Dora Lida Suarez Castro (Uniamazonia)
		\end{tabular} \\ \hline
		Descripci\'on & \begin{tabular}[c]{@{}l@{}}
			El sistema deber\'a suministrar la informaci\'on\\
			respecto a cada uno de los grupos de Investigaci\'on
		\end{tabular} \\ \hline
		SubObjetivos & \\ \hline
		Importancia  & Vital \\ \hline
		Urgencia     & Inmediatamente \\ \hline
		Estado       & En Construcci\'on \\ \hline
		Estabilidad  & Alta \\ \hline
		Comentarios  & \\ \hline
	\end{tabular}
\end{table}


\begin{table}[H]
	\centering
	\caption{Requisito - Gestionar informaci\'on de semilleros de investigaci\'on}
	\label{label_req_gisi}
	\begin{tabular}{|l|l|}
		\hline
		OBJ-005      & \begin{tabular}[c]{@{}l@{}}
			Gestionar informaci\'on de semilleros de investigaci\'on
		\end{tabular} \\ \hline
		Versi\'on    & 001 (2016-08-03) \\ \hline
		Autores      & \begin{tabular}[c]{@{}l@{}}
			Angie Zuleta Cardona (Universidad de la Amazonia)\\
			Derly Viviana  Murcia Serrano (Universidad de la Amazonia) \\ 
			Mateo Ceballos Bermudez (Universidad de la Amazonia)\\ 
			Nicol Dayana Endo Ruiz (Universidad de la Amazonia)
		\end{tabular} \\ \hline
		Fuentes      & \begin{tabular}[c]{@{}l@{}}
			Alberto Fajardo Oliveros (Uniamazonia) \\
			Dora Lida Suarez Castro (Uniamazonia)  \\
		\end{tabular} \\ \hline
		Descripci\'on & \begin{tabular}[c]{@{}l@{}}
			El sistema deber\'a permitir la gesti\'on de la\\
			informaci\'on que corresponde a los semilleros de \\ investigaci\'on.
		\end{tabular} \\  \hline
		SubObjetivos & \\ \hline
		Importancia  & Vital \\ \hline
		Urgencia     & Inmediatamente \\ \hline
		Estado       & En Construcci\'on \\ \hline
		Estabilidad  & Alta \\ \hline
		Comentarios  & Ninguno. \\ \hline
	\end{tabular}
\end{table}

\begin{table}[H]
	\centering
	\caption{Requisito - Gestionar investigador}
	\label{label_req_ginv}
	\begin{tabular}{|l|l|}
		\hline
		OBJ-007      & Gestionar investigador \\ \hline
		Versi\'on    & 002 (2016-08-03) \\ \hline
		Autores      & \begin{tabular}[c]{@{}l@{}}
			Angie Zuleta Cardona (Universidad de la Amazonia)\\   
			Derly Viviana  Murcia Serrano (Universidad de la Amazonia)\\
			Mateo Ceballos Bermudez (Universidad de la Amazonia)\\
			Nicol Dayana Endo Ruiz (Universidad de la Amazonia)
		\end{tabular} \\ \hline
		Fuentes      & \begin{tabular}[c]{@{}l@{}}
			Alberto Fajardo Oliveros (Uniamazonia)\\
			Dora Lida Suarez Castro (Uniamazonia)
		\end{tabular} \\ \hline
		Descripci\'on  & \begin{tabular}[c]{@{}l@{}}
			El sistema deber\'a permitir la gesti\'on de la\\   
			informaci\'on que respecta a los investigadores.
		\end{tabular} \\ \hline
		SubObjetivos & \\ \hline
		Importancia  & Vital \\ \hline
		Urgencia     & Inmediatamente \\ \hline
		Estado       & En Construcci\'on \\ \hline
		Estabilidad  & Alta \\ \hline
		Comentarios  & Ninguna. \\ \hline
	\end{tabular}
\end{table}

\begin{table}[H]
	\centering
	\caption{Requisito - Gestionar par}
	\label{label_req_gpar}
	\begin{tabular}{|l|l|}
		\hline
		OBJ-008      & Gestionar par \\ \hline
		Versi\'on    & 001 (2016-09-09) \\ \hline
		Autores      & \begin{tabular}[c]{@{}l@{}}
			Angie Zuleta Cardona (Universidad de la Amazonia)\\
			Derly Viviana  Murcia Serrano (Universidad de la Amazonia)\\
			Nicol Dayana Endo Ruiz (Universidad de la Amazonia)
		\end{tabular} \\ \hline
		Fuentes      & \begin{tabular}[c]{@{}l@{}}
			Alberto Fajardo Oliveros (Uniamazonia) \\
			Dora Lida Suarez Castro (Uniamazonia)
		\end{tabular} \\ \hline
		Descripcion  & \begin{tabular}[c]{@{}l@{}}
			El sistema deber\'a permitir la gesti\'on de la \\
			informaci\'on referente a los pares evaluadores.
		\end{tabular} \\ \hline
		SubObjetivos & \\ \hline
		Importancia  & Vital \\ \hline
		Urgencia     & Inmediatamente \\ \hline
		Estado       & En Construcci\'on \\ \hline
		Estabilidad  & Alta \\ \hline
		Comentarios  & Ninguno \\ \hline
	\end{tabular}
\end{table}

\begin{table}[H]
	\centering
	\caption{Requisito - Gestionar evaluaci\'on}
	\label{label_req_ge}
	\begin{tabular}{|l|l|}
		\hline
		OBJ-009      & Gestionar evaluaci\'on \\ \hline
		Versi\'on    & 001 (2016-09-09) \\ \hline
		Autores      & \begin{tabular}[c]{@{}l@{}}
			Angie Zuleta Cardona (Universidad de la Amazonia)\\
			Derly Viviana  Murcia Serrano (Universidad de la Amazonia) \\
			Nicol Dayana Endo Ruiz (Universidad de la Amazonia)
		\end{tabular} \\ \hline
		Fuentes      & \begin{tabular}[c]{@{}l@{}}
			Alberto Fajardo Oliveros (Uniamazonia) \\
			Dora Lida Suarez Castro (Uniamazonia)
		\end{tabular} \\ \hline
		Descripci\'on  & \begin{tabular}[c]{@{}l@{}}
			El sistema permitir\'a gestionar la informaci\'on\\
			referente a la evaluaci\'on de las propuestas de \\
			investigaci\'on por parte de los pares evaluadores.
		\end{tabular}\\ \hline
		SubObjetivos & \begin{tabular}[c]{@{}l@{}}
			OBJ-008 Gestionar par
		\end{tabular} \\ \hline
		Importancia  & Vital \\ \hline
		Urgencia     & Inmediatamente \\ \hline
		Estado       & En Construcci\'on \\ \hline
		Estabilidad  & Alta \\ \hline
		Comentarios  & Ninguno \\  \hline
	\end{tabular}
\end{table}

\begin{table}[H]
	\centering
	\caption{Requisito - Gestionar Banco de Elegibles}
	\label{label_req_gbde}
	\begin{tabular}{|l|l|}
		\hline
		OBJ-010      & \begin{tabular}[c]{@{}l@{}}
			Gestionar Banco de Elegibles
		\end{tabular} \\ \hline
		Versi\'on      & 001 (2016-09-09) \\ \hline
		Autores      & \begin{tabular}[c]{@{}l@{}}
			Angie Zuleta Cardona (Universidad de la Amazonia)\\
			Derly Viviana  Murcia Serrano (Universidad de la Amazonia)\\
			Nicol Dayana Endo Ruiz (Universidad de la Amazonia)
		\end{tabular} \\ \hline
		Fuentes      & \begin{tabular}[c]{@{}l@{}}
			Alberto Fajardo Oliveros (Uniamazonia) \\
			Dora Lida Suarez Castro (Uniamazonia)
		\end{tabular} \\ \hline
		Descripci\'on  & \begin{tabular}[c]{@{}l@{}}
			El Sistema deber\'a permitir a los usuarios la\\   gesti\'on del banco de elegibles de propuestas de investigaci\'on, \\
			sea un banco\\   de elegible preliminar o definitivo.
		\end{tabular}        \\ \hline
		SubObjetivos & \begin{tabular}[c]{@{}l@{}}
			OBJ-009 - Gestionar evaluaci\'on
		\end{tabular} \\ \hline
		Importancia  & Vital \\ \hline
		Urgencia     & Inmediatamente \\ \hline
		Estado       & En Construcci\'on \\ \hline
		Estabilidad  & Alta \\ \hline
		Comentarios  & Ninguno. \\ \hline
	\end{tabular}
\end{table}

\begin{table}[H]
	\centering
	\caption{Requisito - Gestionar Preevaluaciones}
	\label{label_req_gp}
	\begin{tabular}{|l|l|}
		\hline
		OBJ-011      & Gestionar Preevaluaciones \\ \hline
		Versi\'on    & 002 (2016-09-09) \\ \hline
		Autores      & \begin{tabular}[c]{@{}l@{}}
			Angie Zuleta Cardona (Universidad de la Amazonia)\\   Derly Viviana  Murcia Serrano\\   (Universidad de la Amazonia) Nicol Dayana Endo Ruiz (Universidad de la\\   Amazonia)
		\end{tabular} \\ \hline
		Fuentes      & \begin{tabular}[c]{@{}l@{}}
			Alberto Fajardo Oliveros (Uniamazonia) \\   
			Dora Lida Suarez Castro (Uniamazonia)
		\end{tabular} \\ \hline
		Descripci\'on  & \begin{tabular}[c]{@{}l@{}}
			El sistema deber\'a permitir la gesti\'on de la informaci\'on\\   referente a las pre- evaluaciones de las propuestas de investigaci\'on
		\end{tabular} \\ \hline
		SubObjetivos & \begin{tabular}[c]{@{}l@{}}
			OBJ-002 - Gestionar propuestas de proyectos
		\end{tabular} \\ \hline
		Importancia  & Vital \\ \hline
		Urgencia     & Inmediatamente \\ \hline
		Estado       & En Construcci\'on \\ \hline
		Estabilidad  & Alta \\ \hline
		Comentarios  & Ninguno \\  \hline
	\end{tabular}
\end{table}

\begin{table}[H]
	\centering
	\caption{Requisito - Gestionar l\'ineas de investigaci\'on}
	\label{label_req_gli}
	\begin{tabular}{|l|l|}
		\hline
		OBJ-012      & \begin{tabular}[c]{@{}l@{}}
			Gestionar L\'ineas de Investigaci\'on
		\end{tabular} \\ \hline
		Versi\'on    & 001 (2016-09-09) \\ \hline
		Autores      & \begin{tabular}[c]{@{}l@{}}
			Angie Zuleta Cardona (Universidad de la Amazonia)\\
			Derly Viviana  Murcia Serrano (Universidad de la Amazonia)\\
			Nicol Dayana Endo Ruiz (Universidad de la Amazonia)
		\end{tabular} \\ \hline
		Fuentes      & \begin{tabular}[c]{@{}l@{}}
			Alberto Fajardo Oliveros (Uniamazonia) \\  
			Dora Lida Suarez Castro (Uniamazonia)
		\end{tabular} \\ \hline
		Descripci\'on & \begin{tabular}[c]{@{}l@{}}
			El Sistema deber\'a permitir la gesti\'on de la\\   informaci\'on que corresponde a las l\'ineas de investigaci\'on relacionadas con\\   los proyectos de investigaci\'on.
		\end{tabular} \\ \hline
		SubObjetivos & \\ \hline
		Importancia  & Vital \\ \hline
		Urgencia     & Inmediatamente \\ \hline
		Estado       & En Construcci\'on  \\ \hline
		Estabilidad  & Alta \\ \hline
		Comentarios  & Ninguno.\\ \hline
	\end{tabular}
\end{table}

\begin{table}[H]
	\centering
	\caption{Requisito - Gestionar inscripci\'on}
	\label{label_req_gins}
	\begin{tabular}{|l|l|}
		\hline
		OBJ-012      & Gestionar inscripci\'on \\ \hline
		Versi\'on    & 001 (2016-09-09) \\ \hline
		Autores      & \begin{tabular}[c]{@{}l@{}}
			Angie Zuleta Cardona (Universidad de la Amazonia)\\ 
			Derly Viviana  Murcia Serrano (Universidad de la Amazonia) \\   
			Nicol Dayana Endo Ruiz (Universidad de la Amazonia)
		\end{tabular} \\ \hline
		Fuentes      & \begin{tabular}[c]{@{}l@{}}
			Alberto Fajardo Oliveros (Uniamazonia) \\
			Dora Lida Suarez Castro (Uniamazonia)
		\end{tabular} \\ \hline
		Descripci\'on & \begin{tabular}[c]{@{}l@{}}
			El Sistema deber\'a permitir la gesti\'on de la\\
			informaci\'on respecto a las distintas inscripciones para propuestas de\\   iniciativa docente en el marco de semilleros y grupos de investigaci\'on
		\end{tabular} \\ \hline
		SubObjetivos & \\ \hline
		Importancia  & Vital \\ \hline
		Urgencia     & Inmediatamente \\ \hline
		Estado       & En Construcci\'on \\ \hline
		Estabilidad  & Alta \\ \hline
		Comentarios  & Ninguno. \\ \hline
	\end{tabular}
\end{table}


\section{Actores del sistema}

\begin{table}[H]
	\centering
	\caption{Actor - Administrador}
	\label{label_actor_admin}
	\begin{tabular}{|l|l|}
		\hline
		ACT-001     & Administrador \\ \hline
		Versi\'on   & 003 (2016-08-04) \\ \hline
		Autores     & \begin{tabular}[c]{@{}l@{}}
			Angie Zuleta Cardona (Universidad de la Amazonia)\\
			Derly Viviana  Murcia Serrano (Universidad de la Amazonia)\\
			Mateo Ceballos Bermudez (Universidad de la Amazonia)\\
			Nicol Dayana Endo Ruiz (Universidad de la Amazonia)
		\end{tabular} \\ \hline
		Fuentes     & \begin{tabular}[c]{@{}l@{}}
			Alberto Fajardo Oliveros (Uniamazonia)\\
			Dora Lida Suarez Castro (Uniamazonia)
		\end{tabular} \\ \hline
		Descripci\'on & \begin{tabular}[c]{@{}l@{}}
			Este actor representa el administrador general\\   del sistema
		\end{tabular} \\ \hline
		Comentarios & Ninguno \\ \hline
	\end{tabular}
\end{table}

\begin{table}[H]
	\centering
	\caption{Actor - Vicerrector}
	\label{label_actor_vice}
	\begin{tabular}{|l|l|}
		\hline
		ACT-002     & Vicerrector \\ \hline
		Versi\'on   & 003 (2016-08-04) \\ \hline
		Autores     & \begin{tabular}[c]{@{}l@{}}
			Angie Zuleta Cardona (Universidad de la Amazonia)\\   
			Derly Viviana  Murcia Serrano (Universidad de la Amazonia)\\
			Mateo Ceballos Bermudez (Universidad de la Amazonia)\\
			Nicol Dayana Endo Ruiz (Universidad de la Amazonia)
		\end{tabular} \\ \hline
		Fuentes     & \begin{tabular}[c]{@{}l@{}}
			Alberto Fajardo Oliveros (Uniamazonia) \\
			Dora Lida Suarez Castro (Uniamazonia)
		\end{tabular}
		\\ \hline
		Descripci\'on & \begin{tabular}[c]{@{}l@{}}
			Este actor representa el vicerrector de\\   investigaciones.\end{tabular} \\ \hline
		Comentarios & Ninguno \\ \hline
	\end{tabular}
\end{table}

\begin{table}[]
	\centering
	\caption{Actor - Auxiliar de investigaci\'on}
	\label{label_actor_auxin}
	\begin{tabular}{|l|l|}
		\hline
		ACT-003     & \begin{tabular}[c]{@{}l@{}}
			Auxiliar de investigaci\'on
		\end{tabular} \\ \hline
		Versi\'on     & 003 (2016-08-04) \\ \hline
		Autores     & \begin{tabular}[c]{@{}l@{}}
			Angie Zuleta Cardona (Universidad de la Amazonia)\\
			Derly Viviana  Murcia Serrano (Universidad de la Amazonia) \\
			Mateo Ceballos Bermudez (Universidad de la Amazonia) \\
			Nicol Dayana Endo Ruiz (Universidad de la Amazonia)
		\end{tabular} \\ \hline
		Fuentes     & \begin{tabular}[c]{@{}l@{}}
			Alberto Fajardo Oliveros (Uniamazonia) \\
			Dora Lida Suarez Castro (Uniamazonia)
		\end{tabular} \\ \hline
		Descripci\'on & \begin{tabular}[c]{@{}l@{}}
			Este actor representa el auxiliar de \\ 
			investigaci\'on que har\'a uso del sistema.
		\end{tabular} \\ \hline
		Comentarios & Ninguno \\ \hline
	\end{tabular}
\end{table}

\begin{table}[H]
	\centering
	\caption{Actor - Investigador}
	\label{label_actor_inv}
	\begin{tabular}{|l|l|}
		\hline
		ACT-004     & Investigador \\ \hline
		Versi\'on   & 003 (2016-08-04) \\ \hline
		Autores     & \begin{tabular}[c]{@{}l@{}}
			Angie Zuleta Cardona (Universidad de la Amazonia)\\
			Derly Viviana  Murcia Serrano (Universidad de la Amazonia) \\
			Mateo Ceballos Bermudez (Universidad de la Amazonia) \\
			Nicol Dayana Endo Ruiz (Universidad de la Amazonia)
		\end{tabular} \\ \hline
		Fuentes     & \begin{tabular}[c]{@{}l@{}}
			Alberto Fajardo Oliveros (Uniamazonia) \\
			Dora Lida Suarez Castro (Uniamazonia)
		\end{tabular} \\ \hline
		Descripci\'on & \begin{tabular}[c]{@{}l@{}}
			Este actor representa el investigador de cada uno\\   de los proyectos de investigaci\'on.
		\end{tabular} \\ \hline
		Comentarios & Ninguno \\ \hline
	\end{tabular}
\end{table}

\begin{table}[]
	\centering
	\caption{Actor - Evaluador}
	\label{label_actor_eval}
	\begin{tabular}{|l|l|}
		\hline
		ACT-005     & Evaluador \\ \hline
		Versi\'on   & 003 (2016-08-04) \\ \hline
		Autores     & \begin{tabular}[c]{@{}l@{}}
			Angie Zuleta Cardona (Universidad de la Amazonia)\\
			Derly Viviana  Murcia Serrano (Universidad de la Amazonia)\\
			Mateo Ceballos Bermudez (Universidad de la Amazonia)\\
			Nicol Dayana Endo Ruiz (Universidad de la Amazonia)
		\end{tabular} \\ \hline
		Fuentes     & \begin{tabular}[c]{@{}l@{}}
			Alberto Fajardo Oliveros (Uniamazonia)\\
			Dora Lida Suarez Castro (Uniamazonia)
		\end{tabular} \\ \hline
		Descripci\'on & \begin{tabular}[c]{@{}l@{}}
			Este actor representa los entes que hacen parte\\   
			del proceso de evaluaci\'on y preevaluaci\'on de las \\
			propuestas de investigaci\'on tanto de los grupos \\
			como de los semilleros. En este caso ser\'ia el Comit\'e de \\  
			investigaciones, Comit\'e de Curr\'iculo, Consejo de Facultad.
		\end{tabular} \\ \hline
		Comentarios & Ninguno \\ \hline
	\end{tabular}
\end{table}

\begin{table}[H]
	\centering
	\caption{Actor - Par evaluador}
	\label{label_actor_pare}
	\begin{tabular}{|l|l|}
		\hline
		ACT-006     & Par evaluador \\ \hline
		Versi\'on     & 003 (2016-08-04) \\ \hline
		Autores     & \begin{tabular}[c]{@{}l@{}}
			Angie Zuleta Cardona (Universidad de la Amazonia)\\   Derly Viviana  Murcia Serrano\\   (Universidad de la Amazonia) Mateo Ceballos\\   Bermudez (Universidad de la Amazonia) Nicol Dayana Endo Ruiz (Universidad de\\   la Amazonia)\end{tabular} \\ \hline
		Fuentes     & \begin{tabular}[c]{@{}l@{}}Alberto \\   Fajardo Oliveros (Uniamazonia) Dora Lida Suarez Castro (Uniamazonia)\end{tabular}                                                                                                                                                     \\ \hline
		Descripci\'on & \begin{tabular}[c]{@{}l@{}}
			Este actor representa el evaluador de cada una de\\   las propuestas de investigaci\'on y de los informes finales de los proyectos.
		\end{tabular} \\ \hline
		Comentarios & Ninguno \\ \hline
	\end{tabular}
\end{table}


\section{Requisitos funcionales}

\begin{table}[H]
	\centering
	\caption{Requisito funcional - RF-001}
	\label{my-label}
	\resizebox{16cm}{!} {	
	\begin{tabular}{|l|l|}
		\hline
		RF-001 & \begin{tabular}[c]{@{}l@{}}
			Registrar convocatoria de propuestas para \\
			proyectos de Grupo de Investigaci\'on
		\end{tabular} \\ \hline
		Versi\'on & 007 (2016-08-04) \\ \hline
		Autores & \begin{tabular}[c]{@{}l@{}}
			Angie Zuleta Cardona (Universidad de la Amazonia)\\   
			Derly Viviana  Murcia Serrano (Universidad de la Amazonia)\\
			Mateo Ceballos Bermudez (Universidad de la Amazonia) \\
			Nicol Dayana Endo Ruiz (Universidad de la Amazonia)
		\end{tabular} \\ \hline
		Fuentes & \begin{tabular}[c]{@{}l@{}}
			Alberto Fajardo Oliveros (Uniamazonia) \\
			Dora Lida Suarez Castro (Uniamazonia)
		\end{tabular} \\ \hline
		Objetivos asociados & \begin{tabular}[c]{@{}l@{}}
			OBJ-001 - Gestionar convocatorias
		\end{tabular} \\ \hline
		Requisitos asociados & \begin{tabular}[c]{@{}l@{}}
			RF-006 \textless Registra propuesta de\\   proyecto\textgreater
		\end{tabular} \\ \hline
		Descripci\'on & \begin{tabular}[c]{@{}l@{}}
			El sistema deber\'a permitir al usuario en \\
			el rol Vicerrector (ACT-002) y Auxiliar de \\
			investigaci\'on (ACT-003) registrar la respectiva\\
			convocatoria  de propuestas de proyectos en \\
			el marco de los Grupos de Investigaci\'on.
		\end{tabular} \\ \hline
		Precondici\'on & \begin{tabular}[c]{@{}l@{}}
			El Usuario debe iniciar sesi\'on en el sistema\\   
			antes de llevar a cabo el proceso de registro.
		\end{tabular} \\ \hline
		Secuencia Normal & Paso \\
		1 & \begin{tabular}[c]{@{}l@{}}
			El Usuario En el rol de Vicerrector (ACT-002) y\\   
			Auxiliar  de investigaci\'on (ACT-003) solicita\\   
			al sistema, iniciar el proceso de registrar \\
			(publicar) la convocatoria.
		\end{tabular} \\
		2 & \begin{tabular}[c]{@{}l@{}}
			El sistema muestra el formulario para el registro\\   de convocatoria.
		\end{tabular} \\
		3 & \begin{tabular}[c]{@{}l@{}}
			El Usuario en el rol de Vicerrector (ACT-002) y\\
			Auxiliar  de investigaci\'on (ACT-003) registra\\   
			(publica) la convocatoria diligenciando el formato \\ 
			teniendo en cuenta los t\'erminos de referencia, que \\ 
			deben ser cumplidos.
		\end{tabular} \\
		4 & \begin{tabular}[c]{@{}l@{}}
			El sistema valida la informaci\'on respectiva de la\\   convocatoria.
		\end{tabular} \\
		5 & \begin{tabular}[c]{@{}l@{}}
			El sistema almacena la informaci\'on de la referida\\   convocatoria en la base de datos.
		\end{tabular} \\
		6 & \begin{tabular}[c]{@{}l@{}}
			El sistema env\'ia un mensaje de confirmaci\'on de\\   
			registro al usuario que est\'e efectuando la operaci\'on.
		\end{tabular} \\
		7 & \begin{tabular}[c]{@{}l@{}}
			El sistema una vez confirme el registro, enviar\'a\\
			una notificaci\'on al correo de los l\'ideres de todos\\
			los grupos de investigaci\'on.
		\end{tabular} \\ \hline
		Postcondici\'on & \begin{tabular}[c]{@{}l@{}}
			Convocatoria registrada con \'exito.
		\end{tabular} \\ \hline
		Excepci\'on & Paso \\
		4 & \begin{tabular}[c]{@{}l@{}}
			El sistema genera una advertencia cuando los datos\\   
			ingresados no son v\'alidos.
		\end{tabular} \\
		5 & \begin{tabular}[c]{@{}l@{}}
			El sistema genera una advertencia cuando \'el \\
			envi\'o no se ha generado de forma exitosa.
		\end{tabular} \\ \hline
		Rendimiento & Paso \begin{tabular}[c]{@{}l@{}}
			Cota de tiempo
		\end{tabular} \\
		2 & 1 segundo \\
		4 & \begin{tabular}[c]{@{}l@{}}
			3 segundos
		\end{tabular} \\
		5 & \begin{tabular}[c]{@{}l@{}}
			3 segundos
		\end{tabular} \\
		Frecuencia Esperada & Semestralmente \\
		Importancia & Vital \\
		Urgencia & Inmediatamente \\
		Estado & En Construcci\'on \\
		Estabilidad & Alta \\
		Comentarios & Ninguno \\
	\end{tabular}
	}
\end{table}

\begin{table}[H]
	\centering
	\caption{My caption}
	\label{my-label}
	\begin{tabular}{|l|l|l|l|l}
		\cline{1-4}
		RF-002 & \begin{tabular}[c]{@{}l@{}}Registrar convocatoria\\   de propuestas para proyectos de Semilleros de investigaci\'on\end{tabular} &  &  &  \\ \cline{1-4}
		Versi\'on & 004 (2016-08-04) &  &  &  \\ \cline{1-4}
		Autores & \begin{tabular}[c]{@{}l@{}}ANGIE ZULETA CARDONA (Universidad de la Amazonia)\\   Derly Viviana  Murcia Serrano\\   (Universidad de la Amazonia) Mateo Ceballos\\   Bermudez (Universidad de la Amazonia) Nicol Dayana Endo Ruiz (Universidad de\\   la Amazonia)\end{tabular} &  &  &  \\ \cline{1-4}
		Fuentes & \begin{tabular}[c]{@{}l@{}}Alberto \\   Fajardo Oliveros (Uniamazonia) Dora Lida Suarez Castro (Uniamazonia)\end{tabular} &  &  &  \\ \cline{1-4}
		Objetivos asociados & \begin{tabular}[c]{@{}l@{}}OBJ-001 \textless Gestionar\\   convocatorias\textgreater\end{tabular} &  &  &  \\ \cline{1-4}
		Requisitos asociados & RF-006 \textless Registra propuesta de proyecto\textgreater &  &  &  \\ \cline{1-4}
		Descripci\'on & \begin{tabular}[c]{@{}l@{}}El sistema deber\'a permitir al usuario en el rol Vicerrector\\   (ACT-002) y Auxiliar  de investigaci\'on\\   (ACT-003)  registrar la respectiva\\   convocatoria  de las propuestas de\\   proyectos en el marco de los Semilleros.\end{tabular} &  &  &  \\ \cline{1-4}
		Precondici\'on & \begin{tabular}[c]{@{}l@{}}El Usuario debe iniciar sesi\'on en el sistema\\   antes de llevar a cabo el proceso de registro.\end{tabular} &  &  &  \\ \cline{1-4}
		\begin{tabular}[c]{@{}l@{}}Secuencia\\   Normal\end{tabular} & \multicolumn{1}{c|}{} & Paso & Acci\'on &  \\ \cline{1-4}
		1 & \begin{tabular}[c]{@{}l@{}}El Usuario En el rol Vicerrector (ACT-002) y\\   Auxiliar  de investigaci\'on\\   (ACT-003)  solicita al sistema, iniciar\\   el proceso de registrar (publicar) convocatoria.\end{tabular} &  &  &  \\ \cline{1-4}
		2 & \begin{tabular}[c]{@{}l@{}}El sistema \\   muestra formulario para registro de convocatoria.\end{tabular} &  &  &  \\ \cline{1-4}
		3 & \begin{tabular}[c]{@{}l@{}}El Usuario en el rol Vicerrector (ACT-002) y\\   Auxiliar  de investigaci\'on\\   (ACT-003)  registra (publica) la\\   convocatoria diligenciando el formato teniendo en cuenta los t\'erminos de\\   referencia de la resoluci\'on, que deben ser cumplidos.\end{tabular} &  &  &  \\ \cline{1-4}
		4 & \begin{tabular}[c]{@{}l@{}}El sistema \\   valida la informaci\'on de la respectiva de la convocatoria.\end{tabular} &  &  &  \\ \cline{1-4}
		5 & \begin{tabular}[c]{@{}l@{}}El sistema almacena la informaci\'on de la referida\\   convocatoria en la base de datos.\end{tabular} &  &  &  \\ \cline{1-4}
		&  & 6 & \begin{tabular}[c]{@{}l@{}}El sistema env�a un mensaje de confirmaci\'on de\\   registro al usuario que est\'e efectuando la operaci\'on.\end{tabular} &  \\ \cline{1-4}
		&  & 7 & \begin{tabular}[c]{@{}l@{}}El sistema una vez confirme el registro, enviar\'a\\   una notificaci\'on al correo de los l�deres de todos los grupos de\\   investigaci\'on.\end{tabular} &  \\ \cline{1-4}
		Postcondici\'on & \begin{tabular}[c]{@{}l@{}}Convocatoria registrada\\   con \'exito.\end{tabular} &  &  &  \\ \cline{1-4}
		Excepci\'on &  & Paso & Acci\'on &  \\ \cline{1-4}
		4 & \begin{tabular}[c]{@{}l@{}}El sistema genera una advertencia cuando los datos\\   ingresados no son v\'alidos.\end{tabular} &  &  &  \\ \cline{1-4}
		&  & 5 & \begin{tabular}[c]{@{}l@{}}El sistema genera una advertencia cuando \'el envi\'o no se\\   ha generado de forma exitosa.\end{tabular} &  \\ \cline{1-4}
		Rendimiento &  & Paso & \begin{tabular}[c]{@{}l@{}}Cota de\\   tiempo\end{tabular} &  \\ \cline{1-4}
		&  & 2 & 1 segundo &  \\ \cline{1-4}
		&  & 4 & \begin{tabular}[c]{@{}l@{}}1\\   segundo\end{tabular} &  \\ \cline{1-4}
		&  & 5 & \begin{tabular}[c]{@{}l@{}}1\\   segundo\end{tabular} &  \\ \cline{1-4}
		Frecuencia Esperada & Semestralmente &  &  &  \\ \cline{1-4}
		Importancia & Vital &  &  &  \\ \cline{1-4}
		Urgencia & Inmediatamente &  &  &  \\ \cline{1-4}
		Estado & En Construcci\'on &  &  &  \\ \cline{1-4}
		Estabilidad & Alta &  &  &  \\ \cline{1-4}
		Comentarios & Ninguno &  &  &  \\ \cline{1-4}
	\end{tabular}
\end{table}

\begin{table}[H]
	\centering
	\caption{My caption}
	\label{my-label}
	\begin{tabular}{|l|l|l|l|l}
		\cline{1-4}
		RF-003 & Consultar convocatoria &  &  &  \\ \cline{1-4}
		Versi\'on & 006 (2016-08-04) &  &  &  \\ \cline{1-4}
		Autores & \begin{tabular}[c]{@{}l@{}}ANGIE ZULETA CARDONA (Universidad de la Amazonia)\\   Derly Viviana  Murcia Serrano\\   (Universidad de la Amazonia) Mateo Ceballos\\   Bermudez (Universidad de la Amazonia) Nicol Dayana Endo Ruiz (Universidad de\\   la Amazonia)\end{tabular} &  &  &  \\ \cline{1-4}
		Fuentes & \begin{tabular}[c]{@{}l@{}}Alberto \\   Fajardo Oliveros (Uniamazonia) Dora Lida Suarez Castro (Uniamazonia)\end{tabular} &  &  &  \\ \cline{1-4}
		Objetivos asociados & \begin{tabular}[c]{@{}l@{}}OBJ-001 \textless Gestionar\\   convocatorias\textgreater\end{tabular} &  &  &  \\ \cline{1-4}
		\begin{tabular}[c]{@{}l@{}}Requisitos\\   asociados\end{tabular} & \begin{tabular}[c]{@{}l@{}}RF-001\\   \textless Registrar convocatoria de propuestas para proyectos de\\   Grupo de Investigaci\'on\textgreater\\   RF-002 \textless Registrar convocatoria de propuestas\\   para Proyectos de Semilleros \textgreater\end{tabular} &  &  &  \\ \cline{1-4}
		Descripci\'on & \begin{tabular}[c]{@{}l@{}}El sistema deber\'a\\   permitir al usuario consultar la informaci\'on de la respectiva convocatoria en\\   el marco de los grupos y semilleros de investigaci\'on.\end{tabular} &  &  &  \\ \cline{1-4}
		Precondici\'on & \begin{tabular}[c]{@{}l@{}}El usuario debe iniciar sesi\'on en el sistema\\   antes de llevar a cabo el proceso de consultar informaci\'on.\\   La convocatoria deber\'a estar previamente\\   registrada.\end{tabular} &  &  &  \\ \cline{1-4}
		\begin{tabular}[c]{@{}l@{}}Secuencia\\   Normal\end{tabular} & \multicolumn{1}{c|}{} & Paso & Acci\'on &  \\ \cline{1-4}
		1 & \begin{tabular}[c]{@{}l@{}}El Usuario solicita al\\   sistema, iniciar el proceso de consulta de una respectiva convocatoria.\end{tabular} &  &  &  \\ \cline{1-4}
		2 & \begin{tabular}[c]{@{}l@{}}El sistema solicita al\\   usuario que identifique si desea consultar Convocatoria en el marco de\\   Semilleros o Grupos de Investigaci\'on.\end{tabular} &  &  &  \\ \cline{1-4}
		3 & \begin{tabular}[c]{@{}l@{}}El sistema muestra la\\   lista de las respectivas convocatorias que han sido registradas previamente.\end{tabular} &  &  &  \\ \cline{1-4}
		4 & \begin{tabular}[c]{@{}l@{}}El sistema solicita al usuario que identifique la convocatoria\\    que desea consultar.\end{tabular} &  &  &  \\ \cline{1-4}
		5 & \begin{tabular}[c]{@{}l@{}}El sistema captura la\\   identificaci\'on anterior.\end{tabular} &  &  &  \\ \cline{1-4}
		6 & \begin{tabular}[c]{@{}l@{}}El sistema verifica en la\\   base de datos y muestra la informaci\'on correspondiente a la convocatoria.\end{tabular} &  &  &  \\ \cline{1-4}
		Postcondici\'on & \begin{tabular}[c]{@{}l@{}}Consulta\\   realizada con \'exito.\end{tabular} &  &  &  \\ \cline{1-4}
		Excepci\'on &  & Paso & Acci\'on &  \\ \cline{1-4}
		6 & \begin{tabular}[c]{@{}l@{}}El\\   sistema genera una advertencia cuando la informaci\'on no es v\'alida.\end{tabular} &  &  &  \\ \cline{1-4}
		Rendimiento &  & Paso & \begin{tabular}[c]{@{}l@{}}Cota de\\   tiempo\end{tabular} &  \\ \cline{1-4}
		2 & \begin{tabular}[c]{@{}l@{}}1\\   segundo\end{tabular} &  &  &  \\ \cline{1-4}
		Frecuencia Esperada & Ninguna &  &  &  \\ \cline{1-4}
		Importancia & Vital &  &  &  \\ \cline{1-4}
		Urgencia & Inmediatamente &  &  &  \\ \cline{1-4}
		Estado & En Construcci\'on &  &  &  \\ \cline{1-4}
		Estabilidad & Alta &  &  &  \\ \cline{1-4}
		Comentarios & Ninguno &  &  &  \\ \cline{1-4}
	\end{tabular}
\end{table}

\begin{table}[H]
	\centering
	\caption{My caption}
	\label{my-label}
	\begin{tabular}{|l|l|l|l|l}
		\cline{1-4} RF-004          & \begin{tabular}[c]{@{}l@{}}
			Consultar informaci\'on de Grupos de Investigaci\'on.
		\end{tabular} & & & \\ \cline{1-4}
		Versi\'on                   & 005 (2016-08-04) & & &  \\ \cline{1-4}
		Autores                     & \begin{tabular}[c]{@{}l@{}}
			ANGIE ZULETA CARDONA (Universidad de la Amazonia)\\
			Derly Viviana  Murcia Serrano(Universidad de la Amazonia)\\
			Mateo Ceballos Bermudez (Universidad de la Amazonia)\\
			Nicol Dayana Endo Ruiz (Universidad de la Amazonia)
		\end{tabular} & & &  \\ \cline{1-4}
		Fuentes                     & \begin{tabular}[c]{@{}l@{}}
			Alberto Fajardo Oliveros (Uniamazonia) \\
			Dora Lida Suarez Castro (Uniamazonia)
		\end{tabular} & & &  \\ \cline{1-4}
		\begin{tabular}[c]{@{}l@{}}
			Objetivos asociados
		\end{tabular} & \begin{tabular}[c]{@{}l@{}}
			OBJ-004 \textless Gestionar informaci\'on de Grupos de\\   Investigaci\'on\textgreater OBJ-005 \textless Gestionar informaci\'on de semilleros de\\   investigaci\'on\textgreater
		\end{tabular} & & &  \\ \cline{1-4}
		Requisitos asociados & & & & \\ \cline{1-4}
		Descripci\'on              & \begin{tabular}[c]{@{}l@{}}
			El sistema deber\'a permitir al usuario consultar\\   
			la informaci\'on de  cada uno de los Grupos de \\ Investigaci\'on.
		\end{tabular} & & &  \\ \cline{1-4}
		Precondici\'on             & \begin{tabular}[c]{@{}l@{}}
			El usuario debe iniciar sesi\'on en el sistema\\
			antes de llevar a cabo el proceso de consultar\\
			informaci\'on de un respectivo grupo de investigaci\'on.\\
			El grupo de investigaci\'on deber\'a estar previamente registrado.
		\end{tabular}              & & &  \\ \cline{1-4}
		\begin{tabular}[c]{@{}l@{}}
			Secuencia Normal
		\end{tabular}    & \multicolumn{1}{c|}{}                                                                                                                                                                                                                                                         & Paso & Acci\'on                                                                                                            &  \\ \cline{1-4}
		1                                                               & \begin{tabular}[c]{@{}l@{}}El Usuario solicita al\\   sistema, iniciar el proceso de consulta de informaci\'on de Grupos de\\   Investigaci\'on.\end{tabular}                                                                                                                     &      &                                                                                                                   &  \\ \cline{1-4}
		2                                                               & \begin{tabular}[c]{@{}l@{}}El sistema despliega la lista de los grupos que\\   se encuentran registrados en el sistema.\end{tabular}                                                                                                                                          &      &                                                                                                                   &  \\ \cline{1-4}
		3                                                               & \begin{tabular}[c]{@{}l@{}}El sistema solicita al usuario que identifique el grupo de investigaci\'on que desea\\   consultar.\end{tabular}                                                                                                                                     &      &                                                                                                                   &  \\ \cline{1-4}
		4                                                               & El sistema captura la identificaci\'on anterior.                                                                                                                                                                                                                                &      &                                                                                                                   &  \\ \cline{1-4}
		5                                                               & \begin{tabular}[c]{@{}l@{}}El sistema verifica en la base de datos y muestra\\   la informaci\'on correspondiente al grupo de investigaci\'on.\end{tabular}                                                                                                                       &      &                                                                                                                   &  \\ \cline{1-4}
		Postcondici\'on                                                   & Consulta realizada con \'exito.                                                                                                                                                                                                                                                 &      &                                                                                                                   &  \\ \cline{1-4}
		Excepci\'on                                                       &                                                                                                                                                                                                                                                                               & Paso & Acci\'on                                                                                                            &  \\ \cline{1-4}
		&                                                                                                                                                                                                                                                                               & 5    & \begin{tabular}[c]{@{}l@{}}El\\   sistema genera una advertencia cuando la informaci\'on no es valida.\end{tabular} &  \\ \cline{1-4}
		Rendimiento                                                     &                                                                                                                                                                                                                                                                               & Paso & \begin{tabular}[c]{@{}l@{}}Cota\\   de tiempo\end{tabular}                                                        &  \\ \cline{1-4}
		&                                                                                                                                                                                                                                                                               & 2    & \begin{tabular}[c]{@{}l@{}}1\\   segundo\end{tabular}                                                             &  \\ \cline{1-4}
		\begin{tabular}[c]{@{}l@{}}Frecuencia\\   Esperada\end{tabular} & Ninguna                                                                                                                                                                                                                                                                       &      &                                                                                                                   &  \\ \cline{1-4}
		Importancia                                                     & Vital                                                                                                                                                                                                                                                                         &      &                                                                                                                   &  \\ \cline{1-4}
		Urgencia                                                        & Inmediatamente                                                                                                                                                                                                                                                                &      &                                                                                                                   &  \\ \cline{1-4}
		Estado                                                          & En Construcci\'on                                                                                                                                                                                                                                                               &      &                                                                                                                   &  \\ \cline{1-4}
		Estabilidad                                                     & Alta                                                                                                                                                                                                                                                                          &      &                                                                                                                   &  \\ \cline{1-4}
		Comentarios                                                     & Ninguno                                                                                                                                                                                                                                                                       &      &                                                                                                                   &  \\ \cline{1-4}
		&                                                                                                                                                                                                                                                                               &      &                                                                                                                   &  \\ \cline{1-4}
		&                                                                                                                                                                                                                                                                               &      &                                                                                                                   &  \\ \cline{1-4}
	\end{tabular}
\end{table}


\begin{table}[]
	\centering
	\caption{My caption}
	\label{my-label}
	\resizebox{\textwidth}{!}{%
		\begin{tabular}{|l|l|l|l|}
			\hline
			RF-005                                                          & \begin{tabular}[c]{@{}l@{}}Consultar\\   informaci\'on de los semilleros de investigaci\'on.\end{tabular}                                                                                                                                                                         &      &                                                                                                                                                             \\ \hline
			Versi\'on                                                         & 004 (2016-08-04)                                                                                                                                                                                                                                                              &      &                                                                                                                                                             \\ \hline
			Autores                                                         & \begin{tabular}[c]{@{}l@{}}ANGIE ZULETA CARDONA (Universidad de la Amazonia)\\   Derly Viviana  Murcia Serrano\\   (Universidad de la Amazonia) Mateo Ceballos\\   Bermudez (Universidad de la Amazonia) Nicol Dayana Endo Ruiz (Universidad de\\   la Amazonia)\end{tabular} &      &                                                                                                                                                             \\ \hline
			Fuentes                                                         & \begin{tabular}[c]{@{}l@{}}Alberto \\   Fajardo Oliveros (Uniamazonia) Dora Lida Suarez Castro (Uniamazonia)\end{tabular}                                                                                                                                                     &      &                                                                                                                                                             \\ \hline
			Objetivos asociados                                             & \begin{tabular}[c]{@{}l@{}}OBJ-005\\   \textless Gestionar informaci\'on de semilleros de investigaci\'on\textgreater\end{tabular}                                                                                                                                                 &      &                                                                                                                                                             \\ \hline
			Requisitos asociados                                            &                                                                                                                                                                                                                                                                               &      &                                                                                                                                                             \\ \hline
			Descripci\'on                                                     & \begin{tabular}[c]{@{}l@{}}El sistema deber\'a permitir consultar la\\   informaci\'on de los semilleros.\end{tabular}                                                                                                                                                            &      &                                                                                                                                                             \\ \hline
			Precondici\'on                                                    & \begin{tabular}[c]{@{}l@{}}El usuario debe iniciar sesi\'on en el sistema\\   antes de llevar a cabo el proceso de consultar informaci\'on.\\   El semillero de\\   investigaci\'on deber\'a estar previamente registrado.\end{tabular}                                               &      &                                                                                                                                                             \\ \hline
			\begin{tabular}[c]{@{}l@{}}Secuencia\\   Normal\end{tabular}    & \multicolumn{1}{c|}{}                                                                                                                                                                                                                                                         & Paso & Acci\'on                                                                                                                                                      \\ \hline
			1                                                               & \begin{tabular}[c]{@{}l@{}}El Usuario \\   solicita al sistema, iniciar el proceso de consulta de informaci\'on de\\   semilleros de investigaci\'on.\end{tabular}                                                                                                                &      &                                                                                                                                                             \\ \hline
			2                                                               & \begin{tabular}[c]{@{}l@{}}El sistema despliega la lista de los semilleros\\   que se encuentran registrados en el sistema.\end{tabular}                                                                                                                                      &      &                                                                                                                                                             \\ \hline
			3                                                               & \begin{tabular}[c]{@{}l@{}}El sistema solicita al usuario que\\   identifique  el semillero de\\   investigaci\'on que desea consultar.\end{tabular}                                                                                                                            &      &                                                                                                                                                             \\ \hline
			4                                                               & El sistema captura la identificaci\'on anterior.                                                                                                                                                                                                                                &      &                                                                                                                                                             \\ \hline
			&                                                                                                                                                                                                                                                                               & 5    & \begin{tabular}[c]{@{}l@{}}El sistema verifica en la base de datos y muestra\\   la informaci\'on correspondiente al semillero de investigaci\'on.\end{tabular} \\ \hline
			Postcondici\'on                                                   & Consulta realizada con \'exito.                                                                                                                                                                                                                                                 &      &                                                                                                                                                             \\ \hline
			Excepci\'on                                                       &                                                                                                                                                                                                                                                                               & Paso & Acci\'on                                                                                                                                                      \\ \hline
			&                                                                                                                                                                                                                                                                               & 5    & \begin{tabular}[c]{@{}l@{}}El\\   sistema genera una advertencia cuando la informaci\'on no es valida.\end{tabular}                                           \\ \hline
			Rendimiento                                                     &                                                                                                                                                                                                                                                                               & Paso & \begin{tabular}[c]{@{}l@{}}Cota\\   de tiempo\end{tabular}                                                                                                  \\ \hline
			&                                                                                                                                                                                                                                                                               & 2    & \begin{tabular}[c]{@{}l@{}}1\\   segundo\end{tabular}                                                                                                       \\ \hline
			&                                                                                                                                                                                                                                                                               & 3    & 3 segundos                                                                                                                                                  \\ \hline
			\begin{tabular}[c]{@{}l@{}}Frecuencia\\   Esperada\end{tabular} & Ninguna                                                                                                                                                                                                                                                                       &      &                                                                                                                                                             \\ \hline
			Importancia                                                     & Vital                                                                                                                                                                                                                                                                         &      &                                                                                                                                                             \\ \hline
			Urgencia                                                        & Inmediatamente                                                                                                                                                                                                                                                                &      &                                                                                                                                                             \\ \hline
			Estado                                                          & En Construcci\'on                                                                                                                                                                                                                                                               &      &                                                                                                                                                             \\ \hline
			Estabilidad                                                     & Alta                                                                                                                                                                                                                                                                          &      &                                                                                                                                                             \\ \hline
			Comentarios                                                     & Ninguno                                                                                                                                                                                                                                                                       &      &                                                                                                                                                             \\ \hline
			&                                                                                                                                                                                                                                                                               &      &                                                                                                                                                             \\ \hline
		\end{tabular}%
	}
\end{table}

% Please add the following required packages to your document preamble:
% \usepackage{graphicx}
\begin{table}[]
	\centering
	\caption{My caption}
	\label{my-label}
	\resizebox{\textwidth}{!}{%
		\begin{tabular}{llll}
			\hline
			\multicolumn{1}{|l|}{RF-006} & \multicolumn{1}{l|}{\begin{tabular}[c]{@{}l@{}}Registrar propuesta de\\   proyecto\end{tabular}} & \multicolumn{1}{l|}{} & \multicolumn{1}{l|}{} \\ \hline
			\multicolumn{1}{|l|}{Versi\'on} & \multicolumn{1}{l|}{006 (2016-08-05)} & \multicolumn{1}{l|}{} & \multicolumn{1}{l|}{} \\ \hline
			\multicolumn{1}{|l|}{Autores} & \multicolumn{1}{l|}{\begin{tabular}[c]{@{}l@{}}ANGIE ZULETA CARDONA (Universidad de la Amazonia) Derly Viviana  Murcia Serrano (Universidad de la Amazonia)\\   Mateo Ceballos Bermudez\\   (Universidad de la Amazonia) Nicol Dayana Endo Ruiz (Universidad de la\\   Amazonia)\end{tabular}} & \multicolumn{1}{l|}{} & \multicolumn{1}{l|}{} \\ \hline
			\multicolumn{1}{|l|}{Fuentes} & \multicolumn{1}{l|}{\begin{tabular}[c]{@{}l@{}}Alberto  Fajardo Oliveros\\   (Uniamazonia) Dora Lida Suarez Castro (Uniamazonia)\end{tabular}} & \multicolumn{1}{l|}{} & \multicolumn{1}{l|}{} \\ \hline
			\multicolumn{1}{|l|}{Objetivos asociados} & \multicolumn{1}{l|}{\begin{tabular}[c]{@{}l@{}}OBJ-002 \textless Gestionar\\   propuestas de proyectos\textgreater\\   OBJ-001 \textless Gestionar\\   convocatorias\textgreater\end{tabular}} & \multicolumn{1}{l|}{} & \multicolumn{1}{l|}{} \\ \hline
			\multicolumn{1}{|l|}{Requisitos asociados} & \multicolumn{1}{l|}{\begin{tabular}[c]{@{}l@{}}RF-001 \textless Registrar\\   convocatoria de propuestas para proyectos de Grupo de Investigaci\'on\textgreater\\   RF-002 \textless Registrar\\   convocatoria de propuestas para proyectos de Semilleros de Investigaci\'on\textgreater\end{tabular}} & \multicolumn{1}{l|}{} & \multicolumn{1}{l|}{} \\ \hline
			\multicolumn{1}{|l|}{Descripci\'on} & \multicolumn{1}{l|}{\begin{tabular}[c]{@{}l@{}}El sistema deber\'a permitir al usuario en el rol de investigador  (ACT-004) registrar su propuesta de\\   proyecto de investigaci\'on en el marco de las convocatorias o ya sea por\\   iniciativa propia.\end{tabular}} & \multicolumn{1}{l|}{} & \multicolumn{1}{l|}{} \\ \hline
			\multicolumn{1}{|l|}{Precondici\'on} & \multicolumn{1}{l|}{\begin{tabular}[c]{@{}l@{}}El usuario debe iniciar sesi\'on en el sistema antes de llevar a cabo el\\   proceso de registrar propuesta de investigaci\'on.\\   La convocatoria deber\'a estar registrada previamente.\end{tabular}} & \multicolumn{1}{l|}{} & \multicolumn{1}{l|}{} \\ \hline
			\multicolumn{1}{|l|}{\begin{tabular}[c]{@{}l@{}}Secuencia\\   Normal\end{tabular}} & \multicolumn{1}{c|}{} & \multicolumn{1}{l|}{Paso} & \multicolumn{1}{l|}{Acci\'on} \\ \hline
			\multicolumn{1}{|l|}{1} & \multicolumn{1}{l|}{\begin{tabular}[c]{@{}l@{}}El Usuario  en el rol de\\   investigador (ACT-004)\\   solicita al sistema iniciar el registro de su propuesta de\\   investigaci\'on.\end{tabular}} & \multicolumn{1}{l|}{} & \multicolumn{1}{l|}{} \\ \hline
			\multicolumn{1}{|l|}{2} & \multicolumn{1}{l|}{\begin{tabular}[c]{@{}l@{}}El sistema solicita al usuario identificar a que convocatoria desea\\   presentar su propuesta.\end{tabular}} & \multicolumn{1}{l|}{} & \multicolumn{1}{l|}{} \\ \hline
			\multicolumn{1}{|l|}{3} & \multicolumn{1}{l|}{\begin{tabular}[c]{@{}l@{}}El sistema muestra el formulario para registrar la propuesta de\\   investigaci\'on\\   \{(FO-M-IV-03-01), (FO-M-IV-03-02), (FO-M-IV-02-05)\}.\end{tabular}} & \multicolumn{1}{l|}{} & \multicolumn{1}{l|}{} \\ \hline
			\multicolumn{1}{|l|}{4} & \multicolumn{1}{l|}{\begin{tabular}[c]{@{}l@{}}El Usuario en el rol de investigador \\   (ACT-004)\\   diligencia el formato de registro de la propuesta de investigaci\'on.\end{tabular}} & \multicolumn{1}{l|}{} & \multicolumn{1}{l|}{} \\ \hline
			\multicolumn{1}{|l|}{5} & \multicolumn{1}{l|}{\begin{tabular}[c]{@{}l@{}}El sistema valida la\\   informaci\'on.\end{tabular}} & \multicolumn{1}{l|}{} & \multicolumn{1}{l|}{} \\ \hline
			\multicolumn{1}{|l|}{6} & \multicolumn{1}{l|}{\begin{tabular}[c]{@{}l@{}}El sistema almacena en la base de datos y env�a un mensaje de\\   confirmaci\'on de registro\end{tabular}} & \multicolumn{1}{l|}{} & \multicolumn{1}{l|}{} \\ \hline
			\multicolumn{1}{|l|}{Postcondici\'on} & \multicolumn{1}{l|}{\begin{tabular}[c]{@{}l@{}}Propuesta registrada\\   con \'exito.\end{tabular}} & \multicolumn{1}{l|}{} & \multicolumn{1}{l|}{} \\ \hline
			\multicolumn{1}{|l|}{Excepci\'on} & \multicolumn{1}{l|}{} & \multicolumn{1}{l|}{Paso} & \multicolumn{1}{l|}{Acci\'on} \\ \hline
			\multicolumn{1}{|l|}{4} & \multicolumn{1}{l|}{\begin{tabular}[c]{@{}l@{}}El sistema genera una advertencia cuando los datos ingresados no son\\   v\'alidos.\end{tabular}} & \multicolumn{1}{l|}{} & \multicolumn{1}{l|}{} \\ \hline
			\multicolumn{1}{|l|}{5} & \multicolumn{1}{l|}{\begin{tabular}[c]{@{}l@{}}El sistema proporciona una advertencia cuando alg�n campo o alg�n\\   requisito este vac�o.\end{tabular}} & \multicolumn{1}{l|}{} & \multicolumn{1}{l|}{} \\ \hline
			\multicolumn{1}{|l|}{Rendimiento} & \multicolumn{1}{l|}{} & \multicolumn{1}{l|}{Paso} & \multicolumn{1}{l|}{\begin{tabular}[c]{@{}l@{}}Cota de\\   tiempo\end{tabular}} \\ \hline
			\multicolumn{1}{|l|}{} & \multicolumn{1}{l|}{} & \multicolumn{1}{l|}{3} & \multicolumn{1}{l|}{\begin{tabular}[c]{@{}l@{}}1\\   Segundo\end{tabular}} \\ \hline
			\multicolumn{1}{|l|}{} & \multicolumn{1}{l|}{} & \multicolumn{1}{l|}{6} & \multicolumn{1}{l|}{\begin{tabular}[c]{@{}l@{}}2\\   Segundo\end{tabular}} \\ \hline
			\multicolumn{1}{|l|}{Frecuencia Esperada} & \multicolumn{1}{l|}{Ninguna.} & \multicolumn{1}{l|}{} & \multicolumn{1}{l|}{} \\ \hline
			\multicolumn{1}{|l|}{Importancia} & \multicolumn{1}{l|}{Vital} & \multicolumn{1}{l|}{} & \multicolumn{1}{l|}{} \\ \hline
			\multicolumn{1}{|l|}{Urgencia} & \multicolumn{1}{l|}{Inmediatamente} & \multicolumn{1}{l|}{} & \multicolumn{1}{l|}{} \\ \hline
			\multicolumn{1}{|l|}{Estado} & \multicolumn{1}{l|}{En Construcci\'on} & \multicolumn{1}{l|}{} & \multicolumn{1}{l|}{} \\ \hline
			\multicolumn{1}{|l|}{Estabilidad} & \multicolumn{1}{l|}{Alta} & \multicolumn{1}{l|}{} & \multicolumn{1}{l|}{} \\ \hline
			Comentarios & Ninguno. &  & 
		\end{tabular}%
	}
\end{table}


% Please add the following required packages to your document preamble:
% \usepackage{graphicx}
\begin{table}[]
	\centering
	\caption{My caption}
	\label{my-label}
	\resizebox{\textwidth}{!}{%
		\begin{tabular}{|l|l|l|l|}
			\hline
			RF-007 & \begin{tabular}[c]{@{}l@{}}Actualizar informaci\'on\\   de propuesta\end{tabular} &  &  \\ \hline
			Versi\'on & 003 (2016-08-05) &  &  \\ \hline
			Autores & \begin{tabular}[c]{@{}l@{}}ANGIE ZULETA CARDONA (Universidad de la Amazonia) Derly Viviana  Murcia Serrano (Universidad de la Amazonia)\\   Mateo Ceballos Bermudez\\   (Universidad de la Amazonia) Nicol Dayana Endo Ruiz (Universidad de la\\   Amazonia)\end{tabular} &  &  \\ \hline
			Fuentes & \begin{tabular}[c]{@{}l@{}}Alberto  Fajardo Oliveros\\   (Uniamazonia) Dora Lida Suarez Castro (Uniamazonia)\end{tabular} &  &  \\ \hline
			Objetivos asociados & \begin{tabular}[c]{@{}l@{}}OBJ-002 \textless Gestionar\\   propuestas de proyectos\textgreater\end{tabular} &  &  \\ \hline
			\begin{tabular}[c]{@{}l@{}}Requisitos\\   asociados\end{tabular} & \begin{tabular}[c]{@{}l@{}}RF-006 \textless Registrar propuesta de proyecto\textgreater RF-008 \textless Consultar\\   propuesta de proyecto \textgreater\end{tabular} &  &  \\ \hline
			Descripci\'on & \begin{tabular}[c]{@{}l@{}}El sistema deber\'a permitir al usuario en el rol de auxiliar de\\   investigaci\'on (ACT-002) o vicerrector (ACT-003) actualizar la informaci\'on de una determinada\\   propuesta de investigaci\'on  si es\\   necesario en el marco de las convocatorias o ya sea por iniciativa propia.\end{tabular} &  &  \\ \hline
			Precondici\'on & \begin{tabular}[c]{@{}l@{}}El usuario el rol de auxiliar de investigaci\'on (ACT-002) o vicerrector (ACT-003) debe iniciar sesi\'on en el sistema antes de\\   llevar a cabo el proceso de actualizar propuesta de investigaci\'on.\\   La propuesta de\\   investigaci\'on deber\'a estar previamente registrada.\end{tabular} &  &  \\ \hline
			\begin{tabular}[c]{@{}l@{}}Secuencia\\   Normal\end{tabular} &  & Paso & Acci\'on \\ \hline
			1 & \begin{tabular}[c]{@{}l@{}}El Usuario en el rol de auxiliar de investigaci\'on (ACT-002) o vicerrector (ACT-003) solicita al sistema iniciar el proceso de\\   actualizar/modificar la informaci\'on de su propuesta de investigaci\'on.\end{tabular} &  &  \\ \hline
			2 & \begin{tabular}[c]{@{}l@{}}El sistema muestra la lista de las propuestas que actualmente se\\   encuentran registradas.\end{tabular} &  &  \\ \hline
			3 & \begin{tabular}[c]{@{}l@{}}El sistema solicita al usuario \\   identificar la propuesta que se desea actualizar.\end{tabular} &  &  \\ \hline
			4 & \begin{tabular}[c]{@{}l@{}}El sistema captura los datos y realiza la consulta en la base de\\   datos.\end{tabular} &  &  \\ \hline
			5 & \begin{tabular}[c]{@{}l@{}}El sistema muestra la interfaz donde se actualizan la informaci\'on de\\   la propuesta identificada anteriormente.\end{tabular} &  &  \\ \hline
			6 & \begin{tabular}[c]{@{}l@{}}El Usuario en el rol de auxiliar de investigaci\'on (ACT-002) o vicerrector (ACT-003) ingresa la informaci\'on necesaria  para la actualizaci\'on de la informaci\'on de\\   su propuesta.\end{tabular} &  &  \\ \hline
			7 & \begin{tabular}[c]{@{}l@{}}El sistema v\'alida la informaci\'on \\   y la almacena en la base de datos.\end{tabular} &  &  \\ \hline
			&  & 8 & El sistema env�a un mensaje de confirmaci\'on de actualizaci\'on. \\ \hline
			Postcondici\'on & \begin{tabular}[c]{@{}l@{}}Propuesta actualizada\\   con \'exito.\end{tabular} &  &  \\ \hline
			Excepci\'on &  & Paso & Acci\'on \\ \hline
			4 & \begin{tabular}[c]{@{}l@{}}El sistema presenta una advertencia cuando la propuesta que se desea\\   actualizar no est\'a registrada.\end{tabular} &  &  \\ \hline
			&  & 7 & El Sistema genera una advertencia cuando la informaci\'on no es v\'alida. \\ \hline
			Rendimiento &  & Paso & \begin{tabular}[c]{@{}l@{}}Cota de\\   tiempo\end{tabular} \\ \hline
			&  & 4 & \begin{tabular}[c]{@{}l@{}}1\\   segundo\end{tabular} \\ \hline
			Frecuencia Esperada & Ninguna &  &  \\ \hline
			Importancia & Vital &  &  \\ \hline
			Urgencia & Inmediatamente &  &  \\ \hline
			Estado & En Construcci\'on &  &  \\ \hline
			Estabilidad & Alta &  &  \\ \hline
			Comentarios & Ninguno &  &  \\ \hline
		\end{tabular}%
	}
\end{table}

% Please add the following required packages to your document preamble:
% \usepackage{graphicx}
\begin{table}[]
	\centering
	\caption{My caption}
	\label{my-label}
	\resizebox{\textwidth}{!}{%
		\begin{tabular}{|l|l|l|l|}
			\hline
			RF-008 & \begin{tabular}[c]{@{}l@{}}Consultar propuesta de\\   proyecto\end{tabular} &  &  \\ \hline
			Versi\'on & 002 (2016-08-06) &  &  \\ \hline
			Autores & \begin{tabular}[c]{@{}l@{}}ANGIE ZULETA CARDONA (Universidad de la Amazonia) Derly Viviana  Murcia Serrano (Universidad de la Amazonia)\\   Mateo Ceballos Bermudez\\   (Universidad de la Amazonia) Nicol Dayana Endo Ruiz (Universidad de la\\   Amazonia)\end{tabular} &  &  \\ \hline
			Fuentes & \begin{tabular}[c]{@{}l@{}}Alberto  Fajardo Oliveros (Uniamazonia)\\   Dora Lida Suarez Castro (Uniamazonia)\end{tabular} &  &  \\ \hline
			Objetivos asociados & \begin{tabular}[c]{@{}l@{}}OBJ-002 \textless Gestionar\\   propuestas de proyectos\textgreater\end{tabular} &  &  \\ \hline
			Requisitos asociados & \begin{tabular}[c]{@{}l@{}}RF-006 \textless Registrar\\   propuesta de proyecto\textgreater\end{tabular} &  &  \\ \hline
			Descripci\'on & \begin{tabular}[c]{@{}l@{}}El sistema deber\'a permitir al usuario en el rol de auxiliar de\\   investigaci\'on (ACT-002),\\   vicerrector (ACT-003),\\   investigador (ACT-004) o evaluador (ACT-005) realizar la consulta de las propuestas de\\   los proyecto de investigaci\'on.\end{tabular} &  &  \\ \hline
			&  &  &  \\ \hline
			Precondici\'on & \begin{tabular}[c]{@{}l@{}}El usuario debe iniciar sesi\'on en el sistema antes de llevar a cabo el\\   proceso de consulta.\\   La propuesta de investigaci\'on deber\'a estar\\   previamente registrada.\end{tabular} &  &  \\ \hline
			\begin{tabular}[c]{@{}l@{}}Secuencia\\   Normal\end{tabular} &  & Paso & Acci\'on \\ \hline
			1 & \begin{tabular}[c]{@{}l@{}}El Usuario  en el rol de\\   auxiliar de investigaci\'on (ACT-002), vicerrector (ACT-003), investigador (ACT-004) o evaluador\\   (ACT-005) solicita al sistema\\   iniciar el proceso de consultar \\   propuesta de investigaci\'on.\end{tabular} &  &  \\ \hline
			2 & \begin{tabular}[c]{@{}l@{}}El sistema despliega la lista de las propuestas que se encuentran\\   registradas en el sistema.\end{tabular} &  &  \\ \hline
			3 & \begin{tabular}[c]{@{}l@{}}El sistema solicita al usuario que identifique  la propuesta de investigaci\'on que desea\\   consultar.\end{tabular} &  &  \\ \hline
			4 & El sistema captura la identificaci\'on anterior. &  &  \\ \hline
			5 & \begin{tabular}[c]{@{}l@{}}El sistema verifica en la base de datos y muestra la informaci\'on\\   correspondiente a la propuesta de investigaci\'on.\end{tabular} &  &  \\ \hline
			Postcondici\'on & \begin{tabular}[c]{@{}l@{}}La propuesta ha sido\\   consultada con \'exito.\end{tabular} &  &  \\ \hline
			Excepci\'on &  & Paso & Acci\'on \\ \hline
			5 & El sistema genera una advertencia cuando la informaci\'on no es v\'alida. &  &  \\ \hline
			Rendimiento &  & Paso & \begin{tabular}[c]{@{}l@{}}Cota de\\   tiempo\end{tabular} \\ \hline
			Frecuencia Esperada & Ninguna. &  &  \\ \hline
			Importancia & Vital &  &  \\ \hline
			Urgencia & Inmediatamente &  &  \\ \hline
			Estado & En Construcci\'on &  &  \\ \hline
			Estabilidad & Alta &  &  \\ \hline
			Comentarios & Ninguno. &  &  \\ \hline
			&  &  &  \\ \hline
			&  &  &  \\ \hline
			&  &  &  \\ \hline
			&  &  &  \\ \hline
		\end{tabular}%
	}
\end{table}

% Please add the following required packages to your document preamble:
% \usepackage{graphicx}
\begin{table}[]
	\centering
	\caption{My caption}
	\label{my-label}
	\resizebox{\textwidth}{!}{%
		\begin{tabular}{|l|l|l|l|}
			\hline
			RF-009 & \begin{tabular}[c]{@{}l@{}}Registrar plan de\\   trabajo\end{tabular} &  &  \\ \hline
			Versi\'on & 002 (2016-08-07) &  &  \\ \hline
			Autores & \begin{tabular}[c]{@{}l@{}}ANGIE ZULETA CARDONA (Universidad de la Amazonia) Derly Viviana  Murcia Serrano (Universidad de la Amazonia)\\   Mateo Ceballos Bermudez\\   (Universidad de la Amazonia) Nicol Dayana Endo Ruiz (Universidad de la\\   Amazonia)\end{tabular} &  &  \\ \hline
			Fuentes & \begin{tabular}[c]{@{}l@{}}Alberto  Fajardo Oliveros\\   (Uniamazonia) Dora Lida Suarez Castro (Uniamazonia)\end{tabular} &  &  \\ \hline
			Objetivos asociados & \begin{tabular}[c]{@{}l@{}}OBJ-003 \textless Gestionar\\   proyectos de Investigaci\'on\textgreater\end{tabular} &  &  \\ \hline
			Requisitos asociados & \begin{tabular}[c]{@{}l@{}}RF-006\\   \textless Registrar propuesta de proyecto\textgreater\end{tabular} &  &  \\ \hline
			Descripci\'on & \begin{tabular}[c]{@{}l@{}}El sistema deber\'a permitir al usuario en el rol de\\   Investigador (ACT-004)  registrar el\\   plan de trabajo del proyecto de investigaci\'on\end{tabular} &  &  \\ \hline
			Precondici\'on & \begin{tabular}[c]{@{}l@{}}El usuario debe iniciar sesi\'on en el sistema antes de\\   llevar a cabo el proceso de registro.\\   La propuesta de investigaci\'on deber\'a estar previamente\\   registrada.\end{tabular} &  &  \\ \hline
			Secuencia Normal &  & Paso & Acci\'on \\ \hline
			1 & \begin{tabular}[c]{@{}l@{}}El Usuario en el rol de investigador (ACT-004)  solicita al sistema, iniciar el proceso de\\   registrar (publicar) el plan de trabajo.\\   de un Proyecto de investigaci\'on.\end{tabular} &  &  \\ \hline
			2 & \begin{tabular}[c]{@{}l@{}}El sistema solicita al\\   usuario identificar a que Proyecto de investigaci\'on desea agregar el plan de\\   trabajo.\end{tabular} &  &  \\ \hline
			3 & \begin{tabular}[c]{@{}l@{}}El sistema muestra formulario para registro del plan de\\   trabajo(PD-M-IV-\\   03), (FO-M-IV-05-03).\end{tabular} &  &  \\ \hline
			4 & \begin{tabular}[c]{@{}l@{}}El Usuario en el rol de investigador (ACT-004)  registra \\   el plan de trabajo del Proyecto diligenciando el formulario\\   correspondiente.\end{tabular} &  &  \\ \hline
			5 & \begin{tabular}[c]{@{}l@{}}El sistema v\'alida la informaci\'on, la almacena en la\\   base de datos y env�a un mensaje de confirmaci\'on de registro.\end{tabular} &  &  \\ \hline
			Postcondici\'on & \begin{tabular}[c]{@{}l@{}}Plan de\\   trabajo registrado con \'exito.\end{tabular} &  &  \\ \hline
			Excepci\'on &  & Paso & Acci\'on \\ \hline
			&  & 5 & \begin{tabular}[c]{@{}l@{}}El sistema genera una\\   advertencia cuando la informaci\'on no es v\'alida.\end{tabular} \\ \hline
			Rendimiento &  & Paso & Cota de tiempo \\ \hline
			&  & 3 & 1 segundo \\ \hline
			Frecuencia Esperada & Ninguna &  &  \\ \hline
			Importancia & Vital &  &  \\ \hline
			Urgencia & Inmediatamente &  &  \\ \hline
			Estado & En Construcci\'on &  &  \\ \hline
			Estabilidad & Alta &  &  \\ \hline
			Comentarios & Ninguno &  &  \\ \hline
		\end{tabular}%
	}
\end{table}

\begin{table}[]
	\centering
	\caption{My caption}
	\label{my-label}
	\resizebox{\textwidth}{!}{ %
		\begin{tabular}{|l|l|l|l|l}
			\cline{1-4}
			RF-010 & Registrar Informes &  &  &  \\ \cline{1-4}
			Versi\'on & 002 (2016-08-06) &  &  &  \\ \cline{1-4}
			Autores & \begin{tabular}[c]{@{}l@{}}ANGIE ZULETA CARDONA (Universidad de la Amazonia) Derly Viviana  Murcia Serrano (Universidad de la Amazonia)\\   Mateo Ceballos Bermudez\\   (Universidad de la Amazonia) Nicol Dayana Endo Ruiz (Universidad de la\\   Amazonia)\end{tabular} &  &  &  \\ \cline{1-4}
			Fuentes & \begin{tabular}[c]{@{}l@{}}Alberto  Fajardo Oliveros\\   (Uniamazonia) Dora Lida Suarez Castro (Uniamazonia)\end{tabular} &  &  &  \\ \cline{1-4}
			Objetivos asociados & \begin{tabular}[c]{@{}l@{}}
				OBJ-006 - Gestionar informes
			\end{tabular} &  &  &  \\ \cline{1-4}
			Requisitos asociados & \begin{tabular}[c]{@{}l@{}}
				RF-006 \textless Registrar propuesta de proyecto
			\end{tabular} &  &  &  \\ \cline{1-4}
			Descripci\'on & \begin{tabular}[c]{@{}l@{}}
				El sistema debe permitir al usuario en el rol de Investigador\\   (ACT-004)  registrar cada uno de los\\   informes sean informes parciales o finales de los proyectos.\end{tabular} &  &  &  \\ \cline{1-4}
			Precondici\'on & \begin{tabular}[c]{@{}l@{}}El usuario debe haber ingresado al sistema para efectuar esta\\   operaci\'on.\\   La propuesta de investigaci\'on deber\'a estar registra previamente.\end{tabular} &  &  &  \\ \cline{1-4}
			\begin{tabular}[c]{@{}l@{}}
				Secuencia Normal
			\end{tabular} &  & Paso & Acci\'on &  \\ \cline{1-4}
			1 & \begin{tabular}[c]{@{}l@{}}
				El Usuario en el rol de investigador solicita al\\   sistema, iniciar el proceso de registrar (publicar) un informe parcial o\\   final de un Proyecto de investigaci\'on.
			\end{tabular} &  &  &  \\ \cline{1-4}
			2 & \begin{tabular}[c]{@{}l@{}}
				El sistema  solicita al usuario identificar a que Proyecto\\   de investigaci\'on desea agregar el informe.
			\end{tabular} &  &  &  \\ \cline{1-4}
			3 & \begin{tabular}[c]{@{}l@{}}
				El sistema solicita al usuario seleccionar que\\   tipo de informe presentar\'a, si final o parcial.
			\end{tabular} &  &  &  \\ \cline{1-4}
			4 & \begin{tabular}[c]{@{}l@{}}
				El sistema muestra formulario para registro de\\   informes GRUPO DE INVESTIGACI�N E INICIATIVA DOCENTE: PARCIAL (FO-M-IV-05-01), FINAL (FO-M-IV-05-02) \\ SEMILLERO DE INVESTIGACI�N (FO-M-IV-04-03).\end{tabular} &  &  &  \\ \cline{1-4}
			5 & \begin{tabular}[c]{@{}l@{}}
				El Usuario en el rol de investigador registra el\\   informe del Proyecto diligenciando el formulario correspondiente.
			\end{tabular} &  &  &  \\ \cline{1-4}
			&  & 6 & \begin{tabular}[c]{@{}l@{}}
				El sistema valida la informaci\'on y la almacena en la base de datos y env\'ia un\\   mensaje de confirmaci\'on de registro.\end{tabular} &  \\ \cline{1-4}
			Postcondici\'on & \begin{tabular}[c]{@{}l@{}}Informe registrado con\\   exito.\end{tabular} &  &  &  \\ \cline{1-4}
			Excepci\'on &  & Paso & Acci\'on &  \\ \cline{1-4}
			6 & \begin{tabular}[c]{@{}l@{}}El\\   sistema genera una advertencia cuando la informaci\'on no es v\'alida.\end{tabular} &  &  &  \\ \cline{1-4}
			Rendimiento &  & Paso & \begin{tabular}[c]{@{}l@{}}Cota de\\   tiempo\end{tabular} &  \\ \cline{1-4}
			&  & 4 & \begin{tabular}[c]{@{}l@{}}1\\   segundo\end{tabular} &  \\ \cline{1-4}
			Frecuencia Esperada & Ninguna. &  &  &  \\ \cline{1-4}
			Importancia & Vital &  &  &  \\ \cline{1-4}
			Urgencia & Inmediatamente &  &  &  \\ \cline{1-4}
			Estado & En Construcci\'on &  &  &  \\ \cline{1-4}
			Estabilidad & Alta &  &  &  \\ \cline{1-4}
			Comentarios & Ninguno. &  &  &  \\ \cline{1-4}
		\end{tabular}
	}
\end{table}


% Please add the following required packages to your document preamble:
% \usepackage{graphicx}
\begin{table}[]
	\centering
	\caption{My caption}
	\label{my-label}
	\resizebox{\textwidth}{!}{%
		\begin{tabular}{|l|l|l|l|}
			\hline
			RF-016 & Asignar par &  &  \\ \hline
			Versi�n & 008 (2016-08-07) &  &  \\ \hline
			Autores & \begin{tabular}[c]{@{}l@{}}ANGIE ZULETA CARDONA (Universidad de la Amazonia)\\   Derly Viviana  Murcia Serrano\\   (Universidad de la Amazonia) Mateo Ceballos\\   Bermudez (Universidad de la Amazonia) Nicol Dayana Endo Ruiz (Universidad de\\   la Amazonia)\end{tabular} &  &  \\ \hline
			Fuentes & \begin{tabular}[c]{@{}l@{}}Alberto \\   Fajardo Oliveros (Uniamazonia) Dora Lida Suarez Castro (Uniamazonia)\end{tabular} &  &  \\ \hline
			Objetivos asociados & \begin{tabular}[c]{@{}l@{}}OBJ-008 \textless Gestionar par\textgreater\\   OBJ-009\textless Gestionar\\   evaluaci�n\textgreater\end{tabular} &  &  \\ \hline
			Requisitos asociados & \begin{tabular}[c]{@{}l@{}}RF-019 \textless Registrar\\   evaluaci�n\textgreater\\   RF-034\\   \textless Registrar par\textgreater\end{tabular} &  &  \\ \hline
			Descripci�n & \begin{tabular}[c]{@{}l@{}}El sistema debe permitir\\   al usuario en el rol de Vicerrector (ACT-002), Auxiliar  de investigaci�n (ACT-003) y Comit� de\\   investigaciones (Evaluador) (ACT-005) efectuar la asignaci�n del par a cada\\   una de las propuestas de investigaci�n o de los informes finales.\\   El usuario deber� tener\\   en cuenta que para asignar un par evaluador a una propuesta estos deber�n\\   coincidir en la l�nea de investigaci�n.\end{tabular} &  &  \\ \hline
			Precondici�n & \begin{tabular}[c]{@{}l@{}}El usuario deber� iniciar\\   sesi�n para llevar a cabo la operaci�n. \\   El par evaluador deber�\\   estar registrado previamente.\\   EN CASO DE PROPUESTAS DE\\   INVESTIGACI�N: La preevaluaci�n de la propuesta de investigaci�n deber� estar\\   registrada y con concepto positivo previamente para realizar la asignaci�n.\\   EN CASO DE INFORMES: El\\   informe del proyecto deber� estar registrado previamente para realizar la\\   asignaci�n.\end{tabular} &  &  \\ \hline
			\begin{tabular}[c]{@{}l@{}}Secuencia\\   Normal\end{tabular} &  & Paso & Acci�n \\ \hline
			1 & \begin{tabular}[c]{@{}l@{}}El Usuario en el rol de\\   Vicerrector (ACT-002), Auxiliar  de\\   investigaci�n (ACT-003) o Comit� de investigaciones en el rol de Evaluador\\   (ACT-005) solicita al sistema, iniciar el proceso de asignar un par evaluador\\   a una propuesta de investigaci�n o un informe.\end{tabular} &  &  \\ \hline
			2 & \begin{tabular}[c]{@{}l@{}}El sistema muestra la lista de las de las propuestas\\   de investigaci�n y los informes que se encuentran registradas actualmente y a\\   las cuales se les va realizar la asignaci�n.\end{tabular} &  &  \\ \hline
			3 & \begin{tabular}[c]{@{}l@{}}El sistema solicita al usuario identificar la\\   propuesta de proyecto de investigaci�n o de informe de proyecto desea asignar\\   el par evaluador teniendo en cuenta los criterios necesarios.\end{tabular} &  &  \\ \hline
			4 & \begin{tabular}[c]{@{}l@{}}El sistema solicita al usuario identifica que\\   tipo de par desea agregar, si es un par externo o interno.\end{tabular} &  &  \\ \hline
			5 & \begin{tabular}[c]{@{}l@{}}El sistema muestra formulario para llevar a cabo\\   la asignaci�n del par a una propuesta de investigaci�n.\end{tabular} &  &  \\ \hline
			6 & \begin{tabular}[c]{@{}l@{}}El Usuario asigna el par a la propuesta de\\   proyecto o al informe final diligenciando el formulario correspondiente.\end{tabular} &  &  \\ \hline
			7 & El sistema valida la informaci�n y la almacena en la base de datos. &  &  \\ \hline
			8 & \begin{tabular}[c]{@{}l@{}}El sistema env�a un mensaje de confirmaci�n de\\   registro.\end{tabular} &  &  \\ \hline
			9 & \begin{tabular}[c]{@{}l@{}}El sistema deber� notificar al correo de la\\   persona que se le ha sido asignado la propuesta o el informe.\end{tabular} &  &  \\ \hline
			Postcondici�n & \begin{tabular}[c]{@{}l@{}}La\\   asignaci�n del par se ha realizado exitosamente.\end{tabular} &  &  \\ \hline
			&  & Paso & Acci�n \\ \hline
			7 & \begin{tabular}[c]{@{}l@{}}El sistema genera una advertencia cuando la\\   informaci�n no es v�lida.\end{tabular} &  &  \\ \hline
			&  & Paso & \begin{tabular}[c]{@{}l@{}}Cota de\\   tiempo\end{tabular} \\ \hline
			2 & \begin{tabular}[c]{@{}l@{}}1\\   segundo\end{tabular} &  &  \\ \hline
			5 & \begin{tabular}[c]{@{}l@{}}1\\   segundo\end{tabular} &  &  \\ \hline
			7 & \begin{tabular}[c]{@{}l@{}}2\\   segundos\end{tabular} &  &  \\ \hline
			Frecuencia Esperada & Semestralmente &  &  \\ \hline
			Importancia & Vital &  &  \\ \hline
			Urgencia & Inmediatamente &  &  \\ \hline
			Estado & En Construcci�n &  &  \\ \hline
			Estabilidad & Alta &  &  \\ \hline
			Comentarios & Ninguno. &  &  \\ \hline
		\end{tabular}%
 	}
\end{table}

% Please add the following required packages to your document preamble:
% \usepackage{graphicx}
\begin{table}[]
	\centering
	\caption{My caption}
	\label{my-label}
	\resizebox{\textwidth}{!}{%
		\begin{tabular}{|l|l|l|l|l}
			\cline{1-4}
			RF-017 & Actualizar par &  &  &  \\ \cline{1-4}
			Versi�n & 006 (2016-08-07) &  &  &  \\ \cline{1-4}
			Autores & \begin{tabular}[c]{@{}l@{}}ANGIE ZULETA CARDONA (Universidad de la Amazonia)\\   Derly Viviana  Murcia Serrano\\   (Universidad de la Amazonia) Nicol Dayana Endo Ruiz (Universidad de la\\   Amazonia)\end{tabular} &  &  &  \\ \cline{1-4}
			Fuentes & \begin{tabular}[c]{@{}l@{}}Alberto \\   Fajardo Oliveros (Uniamazonia) Dora Lida Suarez Castro (Uniamazonia)\end{tabular} &  &  &  \\ \cline{1-4}
			Objetivos asociados & \begin{tabular}[c]{@{}l@{}}OBJ-008 \textless Gestionar par\textgreater\\   OBJ-009 \textless Gestionar\\   evaluaci�n\textgreater\end{tabular} &  &  &  \\ \cline{1-4}
			Requisitos asociados & \begin{tabular}[c]{@{}l@{}}RF-016 \textless Asignar\\   par\textgreater\\   RF-019 \textless Registrar\\   evaluaci�n\textgreater\\   RF-034\\   \textless Registrar par\textgreater\end{tabular} &  &  &  \\ \cline{1-4}
			Descripci�n & \begin{tabular}[c]{@{}l@{}}El sistema debe permitir\\   al usuario en el rol de Vicerrector (ACT-002), Auxiliar  de investigaci�n (ACT-003) y el Comit� de\\   investigaciones en el rol de Evaluador (ACT-005) la actualizaci�n de la\\   informaci�n de los pares que realizan las evaluaciones de las propuestas de\\   investigaci�n o de los informes de los proyectos.\end{tabular} &  &  &  \\ \cline{1-4}
			Precondici�n & \begin{tabular}[c]{@{}l@{}}El usuario deber� iniciar\\   sesi�n para llevar a cabo la operaci�n.  \\   \\   El par evaluador deber�\\   estar registrado previamente.\end{tabular} &  &  &  \\ \cline{1-4}
			\begin{tabular}[c]{@{}l@{}}Secuencia\\   Normal\end{tabular} &  & Paso & Acci�n &  \\ \cline{1-4}
			1 & \begin{tabular}[c]{@{}l@{}}El Usuario en el rol de\\   Vicerrector (ACT-002), Auxiliar  de\\   investigaci�n (ACT-003) y Comit� de Investigaci�n en el rol de Evaluador\\   (ACT-005) solicita al sistema iniciar el proceso de actualizar el par\\   evaluador.\end{tabular} &  &  &  \\ \cline{1-4}
			2 & \begin{tabular}[c]{@{}l@{}}El sistema muestra la lista\\   de los pares que actualmente est�n creados en el sistema.\end{tabular} &  &  &  \\ \cline{1-4}
			3 & \begin{tabular}[c]{@{}l@{}}El sistema solicita al\\   usuario  identificar el par que se\\   desea actualizar.\end{tabular} &  &  &  \\ \cline{1-4}
			4 & \begin{tabular}[c]{@{}l@{}}El sistema captura los\\   datos y realiza la consulta en la base de datos.\end{tabular} &  &  &  \\ \cline{1-4}
			5 & \begin{tabular}[c]{@{}l@{}}El sistema muestra la\\   interfaz donde se actualiza la informaci�n del par evaluador de cada\\   propuesta de Proyecto.\end{tabular} &  &  &  \\ \cline{1-4}
			6 & \begin{tabular}[c]{@{}l@{}}El Usuario en el rol de\\   Vicerrector (ACT-002), Auxiliar  de\\   investigaci�n (ACT-003) o Comit� de investigaciones en el rol de Evaluador\\   (ACT-005) ingresa la informaci�n necesaria para la actualizaci�n del par\\   evaluador.\end{tabular} &  &  &  \\ \cline{1-4}
			7 & \begin{tabular}[c]{@{}l@{}}El sistema v�lida la\\   informaci�n  y la almacena en la base\\   de datos.\end{tabular} &  &  &  \\ \cline{1-4}
			8 & \begin{tabular}[c]{@{}l@{}}El sistema env�a un\\   mensaje de confirmaci�n de actualizaci�n.\end{tabular} &  &  &  \\ \cline{1-4}
			Postcondici�n & \begin{tabular}[c]{@{}l@{}}La\\   actualizaci�n del par se ha realizado con �xito.\end{tabular} &  &  &  \\ \cline{1-4}
			&  & Paso & Acci�n &  \\ \cline{1-4}
			7 & \begin{tabular}[c]{@{}l@{}}El\\   Sistema genera una advertencia cuando la informaci�n no es v�lida.\end{tabular} &  &  &  \\ \cline{1-4}
			Rendimiento &  & Paso & \begin{tabular}[c]{@{}l@{}}Cota de\\   tiempo\end{tabular} &  \\ \cline{1-4}
			&  &  &  &  \\ \cline{1-4}
			Frecuencia Esperada & Ninguna &  &  &  \\ \cline{1-4}
			Importancia & Vital &  &  &  \\ \cline{1-4}
			Urgencia & Inmediatamente &  &  &  \\ \cline{1-4}
			Estado & En Construcci�n &  &  &  \\ \cline{1-4}
			Estabilidad & Alta &  &  &  \\ \cline{1-4}
			Comentarios & Ninguno. &  &  &  \\ \cline{1-4}
		\end{tabular}%
	}
\end{table}

% Please add the following required packages to your document preamble:
% \usepackage{graphicx}
\begin{table}[]
	\centering
	\caption{My caption}
	\label{my-label}
	\resizebox{\textwidth}{!}{%
		\begin{tabular}{|l|l|l|l|l}
			\cline{1-4}
			RF-018 & Consultar par &  &  &  \\ \cline{1-4}
			Versi�n & 004 (2016-08-07) &  &  &  \\ \cline{1-4}
			Autores & \begin{tabular}[c]{@{}l@{}}ANGIE ZULETA CARDONA (Universidad de la Amazonia)\\   Derly Viviana  Murcia Serrano\\   (Universidad de la Amazonia) Nicol Dayana Endo Ruiz (Universidad de la\\   Amazonia)\end{tabular} &  &  &  \\ \cline{1-4}
			Fuentes & \begin{tabular}[c]{@{}l@{}}Alberto \\   Fajardo Oliveros (Uniamazonia) Dora Lida Suarez Castro (Uniamazonia)\end{tabular} &  &  &  \\ \cline{1-4}
			Objetivos asociados & \begin{tabular}[c]{@{}l@{}}OBJ-008\\   \textless Gestionar par\textgreater\\   OBJ-009\\   \textless Gestionar evaluaci�n\textgreater\end{tabular} &  &  &  \\ \cline{1-4}
			Requisitos asociados & \begin{tabular}[c]{@{}l@{}}RF-016 \textless Asignar\\   par\textgreater\\   RF-019 \textless Registrar\\   evaluaci�n\textgreater\\   RF-034\\   \textless Registrar par\textgreater\end{tabular} &  &  &  \\ \cline{1-4}
			Descripci�n & \begin{tabular}[c]{@{}l@{}}El sistema deber�\\   permitir a los usuarios en el rol de Vicerrector (ACT-002), Auxiliar  de investigaci�n (ACT-003) y el Comit� de\\   investigaciones en el rol de Evaluador (ACT-005) realizar la consulta de la\\   informaci�n de cada uno de los pares evaluadores.\end{tabular} &  &  &  \\ \cline{1-4}
			Recondici�n & \begin{tabular}[c]{@{}l@{}}El usuario deber� iniciar sesi�n antes de realizar esta\\   operaci�n. \\   El par evaluador deber�\\   estar registrado.\end{tabular} &  &  &  \\ \cline{1-4}
			\begin{tabular}[c]{@{}l@{}}Secuencia\\   Normal\end{tabular} &  & Paso & Acci�n &  \\ \cline{1-4}
			1 & \begin{tabular}[c]{@{}l@{}}El Usuario en el rol de\\   Vicerrector (ACT-002), Auxiliar  de\\   investigaci�n (ACT-003) o el Comit� de investigaciones en el rol de Evaluador\\   (ACT-005)  solicita al sistema, iniciar\\   el proceso de consulta de informaci�n\\   de los pares evaluadores.\end{tabular} &  &  &  \\ \cline{1-4}
			2 & \begin{tabular}[c]{@{}l@{}}El sistema muestra la\\   lista de los pares que han sido registrados previamente para la evaluaci�n de\\   las propuestas de investigaci�n e informes.\end{tabular} &  &  &  \\ \cline{1-4}
			3 & \begin{tabular}[c]{@{}l@{}}En caso de que se quiera\\   ver la informaci�n m�s a detalle el usuario le indica a el sistema.\end{tabular} &  &  &  \\ \cline{1-4}
			4 & \begin{tabular}[c]{@{}l@{}}El sistema hace visible\\   la informaci�n del par evaluador.\end{tabular} &  &  &  \\ \cline{1-4}
			Postcondici�n & \begin{tabular}[c]{@{}l@{}}Consulta realizada con\\   �xito.\end{tabular} &  &  &  \\ \cline{1-4}
			Excepci�n &  & Paso & Acci�n &  \\ \cline{1-4}
			Rendimiento &  & Paso & \begin{tabular}[c]{@{}l@{}}Cota de\\   tiempo\end{tabular} &  \\ \cline{1-4}
			& 2 & 2 segundos &  &  \\ \cline{1-4}
			&  & 4 & \begin{tabular}[c]{@{}l@{}}1\\   segundo\end{tabular} &  \\ \cline{1-4}
			Frecuencia Esperada & Semestralmente &  &  &  \\ \cline{1-4}
			Importancia & Vital &  &  &  \\ \cline{1-4}
			Urgencia & Inmediatamente &  &  &  \\ \cline{1-4}
			Estado & En Construcci�n &  &  &  \\ \cline{1-4}
			Estabilidad & Alta &  &  &  \\ \cline{1-4}
			Comentarios & Ninguno. &  &  &  \\ \cline{1-4}
		\end{tabular}%
	}
\end{table}

% Please add the following required packages to your document preamble:
% \usepackage{graphicx}
\begin{table}[]
	\centering
	\caption{My caption}
	\label{my-label}
	\resizebox{\textwidth}{!}{%
		\begin{tabular}{|l|l|l|l|l}
			\cline{1-4}
			RF-019 & Registrar evaluaci�n &  &  &  \\ \cline{1-4}
			Versi�n & 001 (2016-09-09) &  &  &  \\ \cline{1-4}
			Autores & \begin{tabular}[c]{@{}l@{}}ANGIE ZULETA CARDONA (Universidad de la Amazonia)\\   Derly Viviana  Murcia Serrano\\   (Universidad de la Amazonia) Nicol Dayana Endo Ruiz (Universidad de la\\   Amazonia)\end{tabular} &  &  &  \\ \cline{1-4}
			Fuentes & \begin{tabular}[c]{@{}l@{}}Alberto \\   Fajardo Oliveros (Uniamazonia) Dora Lida Suarez Castro (Uniamazonia)\end{tabular} &  &  &  \\ \cline{1-4}
			Objetivos asociados & \begin{tabular}[c]{@{}l@{}}OBJ-009 \textless Gestionar\\   evaluaci�n\textgreater\end{tabular} &  &  &  \\ \cline{1-4}
			Requisitos asociados & \begin{tabular}[c]{@{}l@{}}RF-016\\   \textless Asignar par\textgreater\end{tabular} &  &  &  \\ \cline{1-4}
			Descripci�n & \begin{tabular}[c]{@{}l@{}}El sistema deber�\\   permitir al usuario en el rol de Evaluador (ACT-005) siendo el Comit� de\\   investigaciones o en rol de par evaluador (ACT-006), el registro de la\\   informaci�n de cada una de las evaluaciones sean de propuestas de investigaci�n\\   o de informes finales.\\    \\   EN CASO DEL PAR\\   EVALUADOR: El par deber� tener en cuenta los criterios requeridos para\\   registrar la evaluaci�n de las propuestas o de los informes.\\    \\   EN CASO DE SER EL C�MITE\\   DE INVESTIGACIONES: El comit� deber� evaluar a partir del concepto que\\   anteriormente haya registrado el par evaluador ya sea para una propuesta o\\   para un informe.\end{tabular} &  &  &  \\ \cline{1-4}
			Precondici�n & \begin{tabular}[c]{@{}l@{}}El\\   usuario debe haber ingresado al sistema para efectuar esta operaci�n.\\    \\    La evaluaci�n de la propuesta se podr�\\   registrar una vez la preevaluaci�n de la misma, haya tenido un concepto\\   positivo y el par evaluador est� registrado y asignado con anterioridad.\end{tabular} &  &  &  \\ \cline{1-4}
			&  &  &  &  \\ \cline{1-4}
			\begin{tabular}[c]{@{}l@{}}Secuencia\\   Normal\end{tabular} &  & Paso & Acci�n &  \\ \cline{1-4}
			&  & 1 & \begin{tabular}[c]{@{}l@{}}El Usuario en el rol de\\   Evaluador (ACT-005) sea el Comit� de investigaciones  o en el rol de par evaluador (ACT-006) solicita\\   al sistema iniciar el registro de la evaluaci�n ya sea de propuestas de\\   investigaci�n o de informes finales.\end{tabular} &  \\ \cline{1-4}
			&  & 2 & \begin{tabular}[c]{@{}l@{}}El sistema solicita al usuario identificar si la evaluaci�n a\\   registrar ser� para una propuesta de investigaci�n o para un informe de tipo final.\end{tabular} &  \\ \cline{1-4}
			&  & 3 & \begin{tabular}[c]{@{}l@{}}El sistema muestra el formulario\\   para el registro de la evaluaci�n de la propuesta o del informe (Para\\   cualquier caso deber� tener en cuenta los criterios establecidos con\\   anterioridad).\end{tabular} &  \\ \cline{1-4}
			&  & 4 & \begin{tabular}[c]{@{}l@{}}El Usuario en el rol de\\   Evaluador (ACT-005) diligencia el formato de la evaluaci�n de la propuesta de\\   investigaci�n o del informe.\end{tabular} &  \\ \cline{1-4}
			&  & 5 & \begin{tabular}[c]{@{}l@{}}El sistema \\   valida los datos.\end{tabular} &  \\ \cline{1-4}
			&  & 6 & \begin{tabular}[c]{@{}l@{}}El sistema almacena en la\\   base de datos y env�a un mensaje de confirmaci�n de registro\end{tabular} &  \\ \cline{1-4}
			&  & 7 & \begin{tabular}[c]{@{}l@{}}El sistema deber�\\   notificar al l�der del proyecto y a los interesados el concepto de los\\   evaluadores sea de las propuestas o de informes finales.\end{tabular} &  \\ \cline{1-4}
			Postcondici�n & \begin{tabular}[c]{@{}l@{}}Evaluaci�n registrada\\   con �xito\end{tabular} &  &  &  \\ \cline{1-4}
			Excepci�n &  & Paso & Acci�n &  \\ \cline{1-4}
			&  & 5 & \begin{tabular}[c]{@{}l@{}}El sistema genera una advertencia cuando los\\   datos ingresados no son v�lidos.\end{tabular} &  \\ \cline{1-4}
			Rendimiento &  & Paso & Cota de tiempo &  \\ \cline{1-4}
			&  & 3 & \begin{tabular}[c]{@{}l@{}}1\\   segundo\end{tabular} &  \\ \cline{1-4}
			Frecuencia Esperada & Ninguna &  &  &  \\ \cline{1-4}
			Importancia & Vital &  &  &  \\ \cline{1-4}
			Urgencia & Inmediatamente &  &  &  \\ \cline{1-4}
			Estado & En Construcci�n &  &  &  \\ \cline{1-4}
			Estabilidad & Alta &  &  &  \\ \cline{1-4}
			Comentarios & Ninguno &  &  &  \\ \cline{1-4}
		\end{tabular}%
	}
\end{table}

% Please add the following required packages to your document preamble:
% \usepackage{graphicx}
\begin{table}[]
	\centering
	\caption{My caption}
	\label{my-label}
	\resizebox{\textwidth}{!}{%
		\begin{tabular}{|l|l|l|l|}
			\hline
			RF-020 & Actualizar evaluaci�n &  &  \\ \hline
			Versi�n & 001 (2016-09-09) &  &  \\ \hline
			Autores & \begin{tabular}[c]{@{}l@{}}ANGIE ZULETA CARDONA (Universidad de la Amazonia)\\   Derly Viviana  Murcia Serrano\\   (Universidad de la Amazonia) Nicol Dayana Endo Ruiz (Universidad de la\\   Amazonia)\end{tabular} &  &  \\ \hline
			Fuentes & \begin{tabular}[c]{@{}l@{}}Alberto \\   Fajardo Oliveros (Uniamazonia) Dora Lida Suarez Castro (Uniamazonia)\end{tabular} &  &  \\ \hline
			Objetivos asociados & \begin{tabular}[c]{@{}l@{}}OBJ-009 \textless Gestionar\\   evaluaci�n\textgreater\end{tabular} &  &  \\ \hline
			Requisitos asociados & \begin{tabular}[c]{@{}l@{}}RF-019 \textless Registrar\\   evaluaci�n \textgreater\\   RF-016 \textless Asignar\\   Par\textgreater\end{tabular} &  &  \\ \hline
			Descripci�n & \begin{tabular}[c]{@{}l@{}}El sistema deber�\\   permitir al usuario en el rol Evaluador (ACT -005), siendo el Comit� de\\   investigaciones o en el rol de par evaluador (ACT-006), actualizar/modificar\\   la evaluaci�n ya sea de la propuesta de investigaci�n o para un informe de\\   tipo final.\\   La actualizaci�n de la\\   propuesta se dar� solamente en caso de que en las aclaraciones del banco de\\   elegibles, el investigador principal de la propuesta de investigaci�n no est�\\   de acuerdo o  desee que se revisen los\\   criterios o el concepto de la evaluaci�n.\end{tabular} &  &  \\ \hline
			Precondici�n & \begin{tabular}[c]{@{}l@{}}El Usuario debe iniciar\\   sesi�n en el sistema antes de llevar a cabo el proceso de actualizar.\\   La evaluaci�n deber�\\   estar registrada previamente.\end{tabular} &  &  \\ \hline
			\begin{tabular}[c]{@{}l@{}}Secuencia\\   Normal\end{tabular} &  & Paso & Acci�n \\ \hline
			1 & \begin{tabular}[c]{@{}l@{}}El Usuario en el rol de\\   Evaluador (ACT-005), sea el Comit� de investigaciones o en el rol de par evaluador\\   (ACT-006) solicita al sistema iniciar el proceso de actualizar/modificar la evaluaci�n\\   ya sean de propuestas de investigaci�n o de informes parciales.\end{tabular} &  &  \\ \hline
			2 & \begin{tabular}[c]{@{}l@{}}El sistema solicita al\\   usuario identificar si la evaluaci�n a modificar/actualizar ser� para una\\   propuesta de investigaci�n o para un informe de tipo final.\end{tabular} &  &  \\ \hline
			3 & \begin{tabular}[c]{@{}l@{}}El sistema solicita al\\   usuario identificar la evaluaci�n que desea actualizar.\end{tabular} &  &  \\ \hline
			4 & \begin{tabular}[c]{@{}l@{}}El sistema muestra el\\   formulario en el cual se actualiza la evaluaci�n sea de propuestas e informes\\   finales.\end{tabular} &  &  \\ \hline
			5 & \begin{tabular}[c]{@{}l@{}}El Usuario ingresa la\\   informaci�n para la actualizaci�n de la evaluaci�n diligenciando el\\   formulario correspondiente.\end{tabular} &  &  \\ \hline
			6 & \begin{tabular}[c]{@{}l@{}}El sistema  valida la informaci�n y la almacena en la\\   base de datos\end{tabular} &  &  \\ \hline
			7 & \begin{tabular}[c]{@{}l@{}}El sistema deber�\\   notificar al l�der del proyecto y a los interesados que la evaluaci�n ha sido\\   actualizada.\end{tabular} &  &  \\ \hline
			Postcondici�n & \begin{tabular}[c]{@{}l@{}}Evaluaci�n registrada\\   con �xito\end{tabular} &  &  \\ \hline
			Excepci�n &  & Paso & Acci�n \\ \hline
			3 & \begin{tabular}[c]{@{}l@{}}El sistema genera una advertencia cuando se desea\\   actualizar una evaluaci�n que no est� registrada previamente.\end{tabular} &  &  \\ \hline
			7 & \begin{tabular}[c]{@{}l@{}}El sistema genera una advertencia cuando los\\   datos ingresados no son v�lidos.\end{tabular} &  &  \\ \hline
			Rendimiento &  & Paso & Cota de tiempo \\ \hline
			3 & \begin{tabular}[c]{@{}l@{}}1\\   segundo\end{tabular} &  &  \\ \hline
			Frecuencia Esperada &  &  &  \\ \hline
			Importancia & Vital &  &  \\ \hline
			Urgencia & Inmediatamente &  &  \\ \hline
			Estado & En Construcci�n &  &  \\ \hline
			Estabilidad & Alta &  &  \\ \hline
			Comentarios & Ninguno &  &  \\ \hline
		\end{tabular}%
	}
\end{table}


% Please add the following required packages to your document preamble:
% \usepackage{graphicx}
\begin{table}[]
	\centering
	\caption{My caption}
	\label{my-label}
	\resizebox{\textwidth}{!}{%
		\begin{tabular}{|l|l|l|l|}
			\hline
			RF-021 & Consultar evaluaci�n &  &  \\ \hline
			Versi�n & 001 (2016-09-10) &  &  \\ \hline
			Autores & \begin{tabular}[c]{@{}l@{}}ANGIE ZULETA CARDONA (Universidad de la Amazonia)\\   Derly Viviana  Murcia Serrano\\   (Universidad de la Amazonia) Nicol Dayana Endo Ruiz (Universidad de la\\   Amazonia)\end{tabular} &  &  \\ \hline
			Fuentes & \begin{tabular}[c]{@{}l@{}}Alberto \\   Fajardo Oliveros (Uniamazonia) Dora Lida Suarez Castro (Uniamazonia)\end{tabular} &  &  \\ \hline
			Objetivos asociados & \begin{tabular}[c]{@{}l@{}}OBJ-009 \textless Gestionar\\   evaluaci�n\textgreater\end{tabular} &  &  \\ \hline
			Requisitos asociados & \begin{tabular}[c]{@{}l@{}}RF-019 \textless Registrar\\   evaluaci�n \textgreater\end{tabular} &  &  \\ \hline
			Descripci�n & \begin{tabular}[c]{@{}l@{}}El sistema deber�\\   permitir al usuario en el rol de Vicerrector (ACT-002), Auxiliar de\\   investigaci�n (ACT-003), Evaluador (ACT-005) sea el Comit� de investigaciones,\\   en el rol de par evaluador (ACT-006) o Investigador (ACT-001) realizar la\\   consulta de las evaluaciones de las propuestas de investigaci�n o de informes\\   finales de proyectos.\\    \\   En el caso del usuario en\\   el rol de investigador, �ste s�lo podr� consultar las evaluaciones de su\\   propuesta de investigaci�n o del informe final de su proyecto.\end{tabular} &  &  \\ \hline
			Precondici�n & \begin{tabular}[c]{@{}l@{}}El Usuario debe iniciar sesi�n en el sistema\\   antes de llevar a cabo el proceso de consultar la evaluaci�n.\\   La evaluaci�n deber� estar registrada\\   previamente.\end{tabular} &  &  \\ \hline
			\begin{tabular}[c]{@{}l@{}}Secuencia\\   Normal\end{tabular} &  & Paso & Acci�n \\ \hline
			1 & \begin{tabular}[c]{@{}l@{}}El Usuario solicita al sistema, iniciar el\\   proceso de consulta de evaluaci�n de las propuestas o de informes.\end{tabular} &  &  \\ \hline
			2 & \begin{tabular}[c]{@{}l@{}}El sistema solicita al usuario que identifique si\\   desea consultar evaluaciones de propuestas o de informes finales.\end{tabular} &  &  \\ \hline
			3 & \begin{tabular}[c]{@{}l@{}}El sistema muestra la lista de las evaluaciones\\   que han sido registradas previamente.\end{tabular} &  &  \\ \hline
			4 & \begin{tabular}[c]{@{}l@{}}El sistema solicita al usuario que identifique la\\   evaluaci�n que se desea consultar.\end{tabular} &  &  \\ \hline
			5 & El sistema captura los datos &  &  \\ \hline
			& \begin{tabular}[c]{@{}l@{}}El sistema hace visible la informaci�n de la\\   evaluaci�n correspondiente.\end{tabular} &  &  \\ \hline
			Postcondici�n & \begin{tabular}[c]{@{}l@{}}Informaci�n\\   de la evaluaci�n consultada satisfactoriamente.\end{tabular} &  &  \\ \hline
			Excepci�n &  & Paso & Acci�n \\ \hline
			Rendimiento &  & Paso & \begin{tabular}[c]{@{}l@{}}Cota de\\   tiempo\end{tabular} \\ \hline
			Frecuencia Esperada & Ninguna &  &  \\ \hline
			Importancia & Vital &  &  \\ \hline
			Urgencia & Inmediatamente &  &  \\ \hline
			Estado & En Construcci�n &  &  \\ \hline
			Estabilidad & Alta &  &  \\ \hline
			Comentarios & Ninguno &  &  \\ \hline
		\end{tabular}%
	}
\end{table}

% Please add the following required packages to your document preamble:
% \usepackage{graphicx}
\begin{table}[]
	\centering
	\caption{My caption}
	\label{my-label}
	\resizebox{\textwidth}{!}{%
		\begin{tabular}{|l|l|l|l|}
			\hline
			RF-022 & \begin{tabular}[c]{@{}l@{}}Registrar banco de\\   elegibles\end{tabular} &  &  \\ \hline
			Versi�n & 001 (2016-09-10) &  &  \\ \hline
			Autores & \begin{tabular}[c]{@{}l@{}}ANGIE ZULETA CARDONA (Universidad de la Amazonia)\\   Derly Viviana  Murcia Serrano\\   (Universidad de la Amazonia) Nicol Dayana Endo Ruiz (Universidad de la\\   Amazonia)\end{tabular} &  &  \\ \hline
			Fuentes & \begin{tabular}[c]{@{}l@{}}Alberto \\   Fajardo Oliveros (Uniamazonia) Dora Lida Suarez Castro (Uniamazonia)\end{tabular} &  &  \\ \hline
			Objetivos asociados & \begin{tabular}[c]{@{}l@{}}OBJ-010\\   \textless Gestionar Banco de Elegibles\textgreater\end{tabular} &  &  \\ \hline
			Requisitos asociados & \begin{tabular}[c]{@{}l@{}}RF-019 \textless Registrar\\   evaluaci�n\textgreater\\   RF-016 \textless Asignar\\   par\textgreater\end{tabular} &  &  \\ \hline
			Descripci�n & \begin{tabular}[c]{@{}l@{}}El sistema deber�\\   permitir al usuario en el rol de Vicerrector (ACT-002), Auxiliar de\\   investigaci�n (ACT-003), Evaluador (ACT-005) siendo el Comit� de investigaciones\\   o en el rol de  par evaluador (ACT-006)\\    registrar el banco de elegibles, sea\\   preliminar o definitivo.\\    \\   El Banco preliminar de\\   elegibles de las propuestas de investigaci�n se publicar�/registrar�, cuando\\   las evaluaciones de las mismas tengan un concepto positivo. A partir de esta\\   publicaci�n preliminar el investigador tendr� un espacio para realizar una\\   solicitud de aclaraciones en caso de que sea necesario.\\   El banco definitivo de\\   elegibles de las propuestas de investigaci�n se publicar�/registrar�, cuando\\   el lapso de solicitudes de aclaraciones del banco preliminar haya pasado y se\\   hayan realizado los ajustes (en caso de que sean necesarios).\end{tabular} &  &  \\ \hline
			Precondici�n & \begin{tabular}[c]{@{}l@{}}El Usuario debe iniciar\\   sesi�n en el sistema antes de llevar a cabo el proceso de registro.\\   Las evaluaciones de las\\   propuestas deber�n estar registradas previamente y con un concepto positivo.\end{tabular} &  &  \\ \hline
			\begin{tabular}[c]{@{}l@{}}Secuencia\\   Normal\end{tabular} &  & Paso & Acci�n \\ \hline
			1 & \begin{tabular}[c]{@{}l@{}}El Usuario (ACT-001) en el rol de Vicerrector (ACT-002), Auxiliar de investigaci�n\\   (ACT-003), Evaluador (ACT-005) siendo el Comit� de investigaciones o en el\\   rol par evaluador (ACT-006) solicita al sistema, iniciar el proceso de\\   registrar (publicar) el banco de elegibles.\end{tabular} &  &  \\ \hline
			2 & \begin{tabular}[c]{@{}l@{}}El sistema le solicita al usuario que defina qu�\\   tipo de banco desea ingresar, el cual puede ser preliminar o definitivo.\end{tabular} &  &  \\ \hline
			3 & \begin{tabular}[c]{@{}l@{}}El sistema muestra formulario para el registro\\   del banco de elegibles. En el caso del preliminar deber� registrar el espacio\\   para las aclaraciones.\end{tabular} &  &  \\ \hline
			4 & El Usuario registra el banco de elegibles. &  &  \\ \hline
			5 & \begin{tabular}[c]{@{}l@{}}El sistema v�lida la informaci�n, la almacena en\\   la base de datos\end{tabular} &  &  \\ \hline
			6 & \begin{tabular}[c]{@{}l@{}}El sistema env�a un mensaje de confirmaci�n de\\   registro.\end{tabular} &  &  \\ \hline
			&  & 7 & \begin{tabular}[c]{@{}l@{}}El sistema deber� notificar a los interesados la\\   publicaci�n del banco de elegibles.\end{tabular} \\ \hline
			Postcondici�n & \begin{tabular}[c]{@{}l@{}}Banco\\   de elegibles registrado con �xito.\end{tabular} &  &  \\ \hline
			Excepci�n &  & Paso & Acci�n \\ \hline
			&  & 5 & \begin{tabular}[c]{@{}l@{}}El sistema sistema genera una advertencia cuando la\\   informaci�n no es v�lida.\end{tabular} \\ \hline
			Rendimiento &  & Paso & \begin{tabular}[c]{@{}l@{}}Cota de\\   tiempo\end{tabular} \\ \hline
			Frecuencia Esperada & Ninguna. &  &  \\ \hline
			Importancia & Vital &  &  \\ \hline
			Urgencia & Inmediatamente &  &  \\ \hline
			Estado & En Construcci�n &  &  \\ \hline
			Estabilidad & Alta &  &  \\ \hline
			Comentarios & Ninguno &  &  \\ \hline
		\end{tabular}%
	}
\end{table}

% Please add the following required packages to your document preamble:
% \usepackage{graphicx}
\begin{table}[]
	\centering
	\caption{My caption}
	\label{my-label}
	\resizebox{\textwidth}{!}{%
		\begin{tabular}{|l|l|l|l|}
			\hline
			RF-023 & \begin{tabular}[c]{@{}l@{}}Actualizar banco de\\   elegibles\end{tabular} &  &  \\ \hline
			Versi�n & 001 (2016-09-10) &  &  \\ \hline
			Autores & \begin{tabular}[c]{@{}l@{}}ANGIE ZULETA CARDONA (Universidad de la Amazonia)\\   Derly Viviana  Murcia Serrano\\   (Universidad de la Amazonia) Nicol Dayana Endo Ruiz (Universidad de la\\   Amazonia)\end{tabular} &  &  \\ \hline
			Fuentes & \begin{tabular}[c]{@{}l@{}}Alberto \\   Fajardo Oliveros (Uniamazonia) Dora Lida Suarez Castro (Uniamazonia)\end{tabular} &  &  \\ \hline
			Objetivos asociados & \begin{tabular}[c]{@{}l@{}}OBJ-010\\   \textless Gestionar Banco de Elegibles\textgreater\end{tabular} &  &  \\ \hline
			Requisitos asociados & \begin{tabular}[c]{@{}l@{}}RF-022\\   \textless Registrar banco de elegibles\textgreater\end{tabular} &  &  \\ \hline
			Descripci�n & \begin{tabular}[c]{@{}l@{}}El sistema \\   permitir� al usuario en el rol de Vicerrector (ACT-002), Auxiliar de\\   investigaci�n (ACT-003), Evaluador (ACT-005) siendo el Comit� de investigaciones\\   o par evaluador   actualizar los banco de elegibles ya sean\\   preliminares o definitivos.\\   La actualizaci�n del banco de elegibles se har�\\   en caso tal de que sea necesario, por ejemplo cuando el banco definitivo\\   reemplaza el preliminar (el preliminar dejar�a de estar publicado).\end{tabular} &  &  \\ \hline
			Precondici�n & \begin{tabular}[c]{@{}l@{}}El\\   usuario deber� iniciar sesi�n antes de llevar a cabo la operaci�n.\\   El\\   banco de elegibles debe estar previamente registrado.\end{tabular} &  &  \\ \hline
			\begin{tabular}[c]{@{}l@{}}Secuencia\\   Normal\end{tabular} &  & Paso & Acci�n \\ \hline
			1 & \begin{tabular}[c]{@{}l@{}}El Usuario en el rol de Vicerrector (ACT-002),\\   Auxiliar de investigaci�n (ACT-003), Evaluador (ACT-005) (sea el Comit� de\\   investigaciones o par evaluador)  solicita al sistema iniciar el proceso de\\   actualizar/modificar el banco de elegibles.\end{tabular} &  &  \\ \hline
			2 & \begin{tabular}[c]{@{}l@{}}El sistema muestra los bancos de elegibles que se\\   han sido registrados con anterioridad.\end{tabular} &  &  \\ \hline
			3 & \begin{tabular}[c]{@{}l@{}}El sistema solicita al usuario que identifique el\\   banco de elegibles que desea actualizar.\end{tabular} &  &  \\ \hline
			4 & \begin{tabular}[c]{@{}l@{}}El sistema muestra la\\   interfaz donde se actualiza la informaci�n del banco de elegibles.\end{tabular} &  &  \\ \hline
			5 & \begin{tabular}[c]{@{}l@{}}El Usuario ingresa la informaci�n necesaria para\\   la actualizaci�n del banco de elegibles.\end{tabular} &  &  \\ \hline
			6 & El sistema v�lida la informaci�n  y la almacena en la base de datos. &  &  \\ \hline
			7 & \begin{tabular}[c]{@{}l@{}}El sistema notificar� a los interesados la\\   actualizaci�n del banco de elegibles.\end{tabular} &  &  \\ \hline
			Postcondici�n & \begin{tabular}[c]{@{}l@{}}Banco\\   de elegibles actualizado con �xito.\end{tabular} &  &  \\ \hline
			Excepci�n &  & Paso & Acci�n \\ \hline
			6 & \begin{tabular}[c]{@{}l@{}}El sistema genera \\   una advertencia cuando la informaci�n no es v�lida.\end{tabular} &  &  \\ \hline
			Rendimiento &  & Paso & \begin{tabular}[c]{@{}l@{}}Cota de\\   tiempo\end{tabular} \\ \hline
			Frecuencia Esperada & Ninguna. &  &  \\ \hline
			Importancia & Vital &  &  \\ \hline
			Urgencia & Inmediatamente &  &  \\ \hline
			Estado & En Construcci�n &  &  \\ \hline
			Estabilidad & Alta &  &  \\ \hline
			Comentarios & Ninguno. &  &  \\ \hline
		\end{tabular}%
	}
\end{table}

% Please add the following required packages to your document preamble:
% \usepackage{graphicx}
\begin{table}[]
	\centering
	\caption{My caption}
	\label{my-label}
	\resizebox{\textwidth}{!}{%
		\begin{tabular}{|l|l|l|l|}
			\hline
			RF-024 & \begin{tabular}[c]{@{}l@{}}Consultar banco de\\   elegibles\end{tabular} &  &  \\ \hline
			Versi�n & 001 (2016-09-10) &  &  \\ \hline
			Autores & \begin{tabular}[c]{@{}l@{}}ANGIE ZULETA CARDONA (Universidad de la Amazonia)\\   Derly Viviana  Murcia Serrano\\   (Universidad de la Amazonia) Nicol Dayana Endo Ruiz (Universidad de la\\   Amazonia)\end{tabular} &  &  \\ \hline
			Fuentes & \begin{tabular}[c]{@{}l@{}}Alberto \\   Fajardo Oliveros (Uniamazonia) Dora Lida Suarez Castro (Uniamazonia)\end{tabular} &  &  \\ \hline
			Objetivos asociados & \begin{tabular}[c]{@{}l@{}}OBJ-010\\   \textless Gestionar Banco de Elegibles\textgreater\end{tabular} &  &  \\ \hline
			Requisitos asociados & \begin{tabular}[c]{@{}l@{}}RF-022\\   \textless Registrar banco de elegibles\textgreater\end{tabular} &  &  \\ \hline
			Descripci�n & \begin{tabular}[c]{@{}l@{}}El sistema debe permitir a los usuarios consultar\\   los bancos de elegibles sean preliminares o definitivos.\end{tabular} &  &  \\ \hline
			Precondici�n & \begin{tabular}[c]{@{}l@{}}El usuario debe iniciar sesi�n antes de llevar a\\   cabo esta operaci�n.\\   El banco de elegibles deber� estar previamente\\   registrado.\end{tabular} &  &  \\ \hline
			\begin{tabular}[c]{@{}l@{}}Secuencia\\   Normal\end{tabular} &  & Paso & Acci�n \\ \hline
			1 & \begin{tabular}[c]{@{}l@{}}El Usuario solicita\\   al sistema, iniciar el proceso de consultar el banco de elegibles.\end{tabular} &  &  \\ \hline
			2 & \begin{tabular}[c]{@{}l@{}}El sistema muestra la lista de los bancos de\\   elegibles que han sido registrados previamente.\end{tabular} &  &  \\ \hline
			3 & \begin{tabular}[c]{@{}l@{}}En caso de que el usuario desee visualizar la\\   informaci�n en detalle le solicita al sistema.\end{tabular} &  &  \\ \hline
			4 & \begin{tabular}[c]{@{}l@{}}El sistema hace visible la informaci�n del banco\\   de elegibles.\end{tabular} &  &  \\ \hline
			Postcondici�n & \begin{tabular}[c]{@{}l@{}}Consulta realizada con\\   �xito.\end{tabular} &  &  \\ \hline
			Excepci�n &  & Paso & Acci�n \\ \hline
			Rendimiento &  & Paso & \begin{tabular}[c]{@{}l@{}}Cota de\\   tiempo\end{tabular} \\ \hline
			Frecuencia Esperada & Ninguna. &  &  \\ \hline
			Importancia & Vital &  &  \\ \hline
			Urgencia & Inmediatamente &  &  \\ \hline
			Estado & En Construcci�n &  &  \\ \hline
			Estabilidad & Alta &  &  \\ \hline
			Comentarios & Ninguno &  &  \\ \hline
		\end{tabular}%
	}
\end{table}


% Please add the following required packages to your document preamble:
% \usepackage{graphicx}
\begin{table}[]
	\centering
	\caption{My caption}
	\label{my-label}
	\resizebox{\textwidth}{!}{%
		\begin{tabular}{|l|l|l|l|}
			\hline
			RF-025 & \begin{tabular}[c]{@{}l@{}}Registrar proyectos de\\   investigaci�n.\end{tabular} &  &  \\ \hline
			Versi�n & 001 (2016-09-10) &  &  \\ \hline
			Autores & \begin{tabular}[c]{@{}l@{}}ANGIE ZULETA CARDONA (Universidad de la Amazonia)\\   Derly Viviana  Murcia Serrano\\   (Universidad de la Amazonia) Nicol Dayana Endo Ruiz (Universidad de la\\   Amazonia)\end{tabular} &  &  \\ \hline
			Fuentes & \begin{tabular}[c]{@{}l@{}}Alberto \\   Fajardo Oliveros (Uniamazonia) Dora Lida Suarez Castro (Uniamazonia)\end{tabular} &  &  \\ \hline
			Objetivos asociados & \begin{tabular}[c]{@{}l@{}}OBJ-003\\   \textless Gestionar proyectos de Investigaci�n\textgreater\end{tabular} &  &  \\ \hline
			Requisitos asociados & \begin{tabular}[c]{@{}l@{}}RF-019\\   \textless Registrar evaluaci�n \textgreater\\   RF-023\\   \textless Registrar banco de elegibles \textgreater\end{tabular} &  &  \\ \hline
			Descripci�n & \begin{tabular}[c]{@{}l@{}}El sistema deber� permitir  al usuario en el rol de Vicerrector\\   (ACT-002), Auxiliar de investigaci�n (ACT-003) e Investigador (ACT-001) registrar\\   la informaci�n de cada una de los proyectos realizadas por los\\   investigadores, siendo estos las propuestas que han sido aprobadas y\\   consolidadas mediante la evaluaci�n que ha realizado el par evaluador y el\\   comit� de investigaciones.\end{tabular} &  &  \\ \hline
			Precondici�n & \begin{tabular}[c]{@{}l@{}}El usuario debe haber ingresado al sistema para\\   efectuar esta operaci�n.\\   La evaluaci�n de la propuesta debi� tener\\   concepto positivo.\\   La propuesta debi� estar en el banco definitivo\\   de elegibles.\end{tabular} &  &  \\ \hline
			\begin{tabular}[c]{@{}l@{}}Secuencia\\   Normal\end{tabular} &  & Paso & Acci�n \\ \hline
			1 & \begin{tabular}[c]{@{}l@{}}El Usuario en el rol de Vicerrector\\   (ACT-002),  Auxiliar de investigaci�n\\   (ACT-003) e Investigador (ACT-001) solicita al sistema, iniciar el proceso de\\   registrar Proyecto de investigaci�n.\end{tabular} &  &  \\ \hline
			2 & \begin{tabular}[c]{@{}l@{}}El sistema muestra formulario para el registro de\\   proyectos de investigaci�n\end{tabular} &  &  \\ \hline
			3 & El Usuario registra el Proyecto de investigaci�n. &  &  \\ \hline
			4 & \begin{tabular}[c]{@{}l@{}}El sistema v�lida la informaci�n, la almacena en\\   la base de datos\end{tabular} &  &  \\ \hline
			&  & 5 & \begin{tabular}[c]{@{}l@{}}El sistema  env�a un mensaje de confirmaci�n de\\   registro.\end{tabular} \\ \hline
			Postcondici�n & \begin{tabular}[c]{@{}l@{}}Proyecto registrado\\   con �xito.\end{tabular} &  &  \\ \hline
			Excepci�n &  & Paso & Acci�n \\ \hline
			&  &  &  \\ \hline
			&  &  &  \\ \hline
			Rendimiento &  & Paso & \begin{tabular}[c]{@{}l@{}}Cota de\\   tiempo\end{tabular} \\ \hline
			Frecuencia Esperada & Ninguna &  &  \\ \hline
			Importancia & Vital &  &  \\ \hline
			Urgencia & Inmediatamente &  &  \\ \hline
			Estado & En Construcci�n &  &  \\ \hline
			Estabilidad & Alta &  &  \\ \hline
			Comentarios & Ninguno &  &  \\ \hline
		\end{tabular}%
	}
\end{table}

% Please add the following required packages to your document preamble:
% \usepackage{graphicx}
\begin{table}[]
	\centering
	\caption{My caption}
	\label{my-label}
	\resizebox{\textwidth}{!}{%
		\begin{tabular}{|l|l|l|l|}
			\hline
			RF-026 & \begin{tabular}[c]{@{}l@{}}Consultar proyecto de\\   investigaci�n\end{tabular} &  &  \\ \hline
			Versi�n & 001 (2016-09-10) &  &  \\ \hline
			Autores & \begin{tabular}[c]{@{}l@{}}ANGIE ZULETA CARDONA (Universidad de la Amazonia)\\   Derly Viviana  Murcia Serrano (Universidad\\   de la Amazonia) Nicol Dayana Endo Ruiz (Universidad de la Amazonia)\end{tabular} &  &  \\ \hline
			Fuentes & \begin{tabular}[c]{@{}l@{}}Alberto \\   Fajardo Oliveros (Uniamazonia) Dora Lida Suarez Castro (Uniamazonia)\end{tabular} &  &  \\ \hline
			Objetivos asociados & \begin{tabular}[c]{@{}l@{}}OBJ-003\\   \textless Gestionar proyectos de Investigaci�n\textgreater\end{tabular} &  &  \\ \hline
			Requisitos asociados & \begin{tabular}[c]{@{}l@{}}RF-025\\   \textless Registrar proyectos de investigaci�n.\textgreater\end{tabular} &  &  \\ \hline
			Descripci�n & \begin{tabular}[c]{@{}l@{}}El sistema deber� permitir consultar la\\   informaci�n de cada uno de los proyectos de investigaci�n.\end{tabular} &  &  \\ \hline
			Precondici�n & \begin{tabular}[c]{@{}l@{}}El\\   usuario deber� iniciar sesi�n antes de realizar esta operaci�n.\end{tabular} &  &  \\ \hline
			\begin{tabular}[c]{@{}l@{}}Secuencia\\   Normal\end{tabular} &  & Paso & Acci�n \\ \hline
			1 & \begin{tabular}[c]{@{}l@{}}El Usuario (ACT-001) solicita al sistema, iniciar el proceso de consulta de\\   proyectos de investigaci�n.\end{tabular} &  &  \\ \hline
			2 & \begin{tabular}[c]{@{}l@{}}El sistema muestra el formulario para llevar a\\   cabo la consulta.\end{tabular} &  &  \\ \hline
			3 & \begin{tabular}[c]{@{}l@{}}El sistema solicita al usuario que identifique el\\   Proyecto de investigaci�n que se desea consultar.\end{tabular} &  &  \\ \hline
			4 & \begin{tabular}[c]{@{}l@{}}El sistema captura los\\   datos, verifica si proyecto est� registrado en la base de datos.\end{tabular} &  &  \\ \hline
			5 & \begin{tabular}[c]{@{}l@{}}El sistema hace visible la informaci�n del\\   proyecto de investigaci�n\end{tabular} &  &  \\ \hline
			Postcondici�n & \begin{tabular}[c]{@{}l@{}}Consulta realizada con\\   �xito.\end{tabular} &  &  \\ \hline
			Excepci�n &  & Paso & Acci�n \\ \hline
			Rendimiento &  & Paso & \begin{tabular}[c]{@{}l@{}}Cota de\\   tiempo\end{tabular} \\ \hline
			Frecuencia Esperada & Ninguna. &  &  \\ \hline
			Importancia & Vital &  &  \\ \hline
			Urgencia & Inmediatamente &  &  \\ \hline
			Estado & En Construcci�n &  &  \\ \hline
			Estabilidad & Alta &  &  \\ \hline
			Comentarios & Ninguno &  &  \\ \hline
		\end{tabular}%
	}
\end{table}

% Please add the following required packages to your document preamble:
% \usepackage{graphicx}
\begin{table}[]
	\centering
	\caption{My caption}
	\label{my-label}
	\resizebox{\textwidth}{!}{%
		\begin{tabular}{|l|l|l|l|}
			\hline
			RF-034 & \begin{tabular}[c]{@{}l@{}}Registrar\\   par\end{tabular} &  &  \\ \hline
			Versi�n & 005 (2016-08-04) &  &  \\ \hline
			Autores & \begin{tabular}[c]{@{}l@{}}ANGIE ZULETA CARDONA (Universidad de la Amazonia)\\   Derly Viviana  Murcia Serrano\\   (Universidad de la Amazonia) Mateo Ceballos\\   Bermudez (Universidad de la Amazonia) Nicol Dayana Endo Ruiz (Universidad de\\   la Amazonia)\end{tabular} &  &  \\ \hline
			Fuentes & \begin{tabular}[c]{@{}l@{}}Alberto  Fajardo\\   Oliveros (Uniamazonia) Dora Lida Suarez Castro (Uniamazonia)\end{tabular} &  &  \\ \hline
			\begin{tabular}[c]{@{}l@{}}Objetivos\\   asociados\end{tabular} & \begin{tabular}[c]{@{}l@{}}OBJ-008 \textless Gestionar\\   par\textgreater\\   OBJ-009 \textless Gestionar\\   evaluaci�n\textgreater\end{tabular} &  &  \\ \hline
			Requisitos asociados & \begin{tabular}[c]{@{}l@{}}RF-019 \textless Registrar\\   evaluaci�n\textgreater\end{tabular} &  &  \\ \hline
			Descripci�n & \begin{tabular}[c]{@{}l@{}}El sistema debe permitir al usuario en el rol de\\   Vicerrector (ACT-002), Auxiliar  de\\   investigaci�n (ACT-003) y Comit� de investigaciones (Evaluador) (ACT-005)\\   efectuar el registro del par quienes realizar�n las evaluaciones de las propuestas\\   de investigaci�n o de los informes sean parciales o finales.\end{tabular} &  &  \\ \hline
			Precondici�n & \begin{tabular}[c]{@{}l@{}}El usuario deber� iniciar sesi�n para llevar a\\   cabo la operaci�n. \\   El par evaluador no debe haberse registrado\\   previamente\end{tabular} &  &  \\ \hline
			\begin{tabular}[c]{@{}l@{}}Secuencia\\   Normal\end{tabular} &  & Paso & Acci�n \\ \hline
			1 & \begin{tabular}[c]{@{}l@{}}El Usuario usuario en el\\   rol de Vicerrector (ACT-002), Auxiliar \\   de investigaci�n (ACT-003) o Comit� de investigaciones (Evaluador)\\   (ACT-005)  solicita al sistema, iniciar\\   el proceso de registro de la informaci�n del par evaluador.\end{tabular} &  &  \\ \hline
			2 & \begin{tabular}[c]{@{}l@{}}El sistema solicita al\\   usuario que identifique si el par a registrar es interno o externo.\end{tabular} &  &  \\ \hline
			3 & \begin{tabular}[c]{@{}l@{}}El sistema muestra el formulario para llevar a\\   cabo el registro.\end{tabular} &  &  \\ \hline
			4 & \begin{tabular}[c]{@{}l@{}}El usuario realiza el registro, diligenciando el\\   formulario correspondiente.\end{tabular} &  &  \\ \hline
			5 & El sistema v�lida la informaci�n &  &  \\ \hline
			6 & \begin{tabular}[c]{@{}l@{}}El sistema almacena la informaci�n en la base de\\   datos\end{tabular} &  &  \\ \hline
			7 & \begin{tabular}[c]{@{}l@{}}El sistema env�a un mensaje de confirmaci�n de\\   registro del par al usuario que est� efectuando la operaci�n.\end{tabular} &  &  \\ \hline
			Postcondici�n & El par ha sido registrado con �xito. &  &  \\ \hline
			Excepci�n &  & Paso & Acci�n \\ \hline
			4 & \begin{tabular}[c]{@{}l@{}}El\\   Sistema genera una advertencia cuando el par se ha registrado previamente.\end{tabular} &  &  \\ \hline
			5 & \begin{tabular}[c]{@{}l@{}}El\\   Sistema genera una advertencia cuando los datos no son v�lidos.\end{tabular} &  &  \\ \hline
			Rendimiento &  & Paso & \begin{tabular}[c]{@{}l@{}}Cota de\\   tiempo\end{tabular} \\ \hline
			Frecuencia Esperada & Ninguna &  &  \\ \hline
			Importancia & Vital &  &  \\ \hline
			Urgencia & Inmediatamente &  &  \\ \hline
			Estado & En Construcci�n &  &  \\ \hline
			Estabilidad & Alta &  &  \\ \hline
			Comentarios & Ninguno &  &  \\ \hline
		\end{tabular}%
	}
\end{table}
	
	% Finalizaci�n
	%
%\makeglossaries
%\loadglsentries{glossaires}
%
%\newglossaryentry{latex}
%{
%	name=latex,
%	description={Is a mark up language specially suited 
%		for scientific documents}
%}
%
%\newglossaryentry{maths}
%{
%	name=mathematics,
%	description={Mathematics is what mathematicians do}
%}

%\title{How to create a glossary}
%\author{ }
%\date{ }

%\maketitle

%The \Gls{latex} typesetting markup language is specially suitable 
%for documents that include \gls{maths}. 

%\clearpage

%\printglossaries
	\section[CONCLUSIONES]{CONCLUSIONES}
	
	\bibliographystyle{alpha}
	\bibliography{referencias}
	
	\section[ANEXOS]{ANEXOS}
\end{document}